%%%%%%%%%%%%%%%%%%%%%%%%%%%%%%%%%%%%%%%%%%%%%%%%%%%%%%%%%%%%%%%%%%%%%%%%%%%%%%%%
% FLAGS
%
\newif\iffull
\fulltrue
%\fullfalse
%%%%%%%%%%%%%%%%%%%%%%%%%%%%%%%%%%%%%%%%%%%%%%%%%%%%%%%%%%%%%%%%%%%%%%%%%%%%%%%%


%%%%%%%%%%%%%%%%%%%%%%%%%%%%%%%%%%%%%%%%%%%%%%%%%%%%%%%%%%%%%%%%%%%%%%%%%%%%%%%
\iffull
  \documentclass[11pt,letterpaper]{article}
  \IfFileExists{dtrt.sty}{%
    \usepackage[notes=true,later=true]{dtrt}%
  }{%
    \usepackage{xcolor}
    \usepackage{hyperref}
  }
  \usepackage{fullpage}
\else
  \documentclass[letter]{llncs}
  \usepackage[notes=xxx,llncssubsub]{dtrt}
\fi
\providecommand{\dtcolornote}[3][]{\textcolor{#2}{[\textbf{#1}: #3]}}
\providecommand{\parhead}[1]{\paragraph{#1}}

% Bibliography entries are included at the end of this standalone note.
\newcommand*{\bibfont}{\small}

%%%%%%%%%%%%%%%%%%%%%%%%%%%%%%%%%%%%%%%%%%%%%%%%%%%%%%%%%%%%%%%%%%%%%%%%%%%%%%%%


%%%%%%%%%%%%%%%%%%%%%%%%%%%%%%%%%%%%%%%%%%%%%%%%%%%%%%%%%%%%%%%%%%%%%%%%%%%%%%%%
% PACKAGES
\iffull
\usepackage{amsthm}
\else
\let\proof\relax
\let\endproof\relax
\usepackage{amsthm}
\fi
\usepackage[T1]{fontenc}
\usepackage{float}
\usepackage{amssymb}
\usepackage{amsmath}
\usepackage{setspace}
\usepackage{bbm}
\usepackage{bm}
\usepackage{paralist}
\usepackage{enumitem}
  \setlist[itemize]{leftmargin=*}
  \setlist[enumerate]{leftmargin=*}
\usepackage{soul}
\usepackage{xspace}


\usepackage[nameinlink]{cleveref}
  \creflabelformat{equation}{#2#1#3}
\usepackage{multirow}
\usepackage{hhline}
\usepackage{makecell}
\usepackage{booktabs}
\usepackage{mathtools}

\usepackage{paralist}

\usepackage{tikz}
\usetikzlibrary{decorations.pathreplacing}
\usetikzlibrary{shapes.geometric}
\usetikzlibrary{arrows}
\usetikzlibrary{patterns}

\usepackage{color, colortbl}
\usepackage{xfrac}

\iffull
\usepackage[margin=4mm,small,labelfont=bf]{caption}
\fi

\usepackage{subcaption}
\usepackage{diagbox}
%%%%%%%%%%%%%%%%%%%%%%%%%%%%%%%%%%%%%%%%%%%%%%%%%%%%%%%%%%%%%%%%%%%%%%%%%%%%%%%%



%%%%%%%%%%%%%%%%%%%%%%%%%%%%%%%%%%%%%%%%%%%%%%%%%%%%%%%%%%%%%%%%%%%%%%%%%%%%%%%%
% AUTHOR NOTES
%
\newcommand{\ale}[1]{\dtcolornote[Ale]{red}{#1}}
\newcommand{\eylon}[1]{\dtcolornote[Eylon]{blue}{#1}}
\newcommand{\gal}[1]{\dtcolornote[Gal]{cyan}{#1}}
\newcommand{\giacomo}[1]{\dtcolornote[Giacomo]{violet}{#1}}
%%%%%%%%%%%%%%%%%%%%%%%%%%%%%%%%%%%%%%%%%%%%%%%%%%%%%%%%%%%%%%%%%%%%%%%%%%%%%%%%



%%%%%%%%%%%%%%%%%%%%%%%%%%%%%%%%%%%%%%%%%%%%%%%%%%%%%%%%%%%%%%%%%%%%%%%%%%%%%%%%
% ENVIRONMENTS
%

% amsthm defs
% non-nested
\newtheorem{itheorem}{Theorem}%[section]
\newtheorem{iconjecture}{Conjecture}%[section]
\newtheorem{idefinition}{Definition}
\newtheorem{ilemma}{Lemma}%[section]
\newtheorem{icorollary}{Corollary}%[section]
\newtheorem{iclaim}{Claim}
\newtheorem{ifact}{Fact}
\iffull
\newtheorem{theorem}{Theorem}[section]
\newtheorem{lemma}[theorem]{Lemma}
\newtheorem{proposition}[theorem]{Proposition}
\newtheorem{corollary}[theorem]{Corollary}
\newtheorem{definition}[theorem]{Definition}
\newtheorem{assumption}[theorem]{Assumption}
\newtheorem{observation}[theorem]{Observation}
\newtheorem{claim}[theorem]{Claim}
\newtheorem{conjecture}[theorem]{Conjecture}

\else
\let\claim\relax % undefine the environment
\spnewtheorem{claim}[lemma]{Claim}{\bfseries}{\rmfamily}
\crefname{claim}{claim}{claims}
\fi


\newtheorem{folklore}[theorem]{Folklore}
\crefname{conjecture}{conjecture}{conjectures}
\newtheorem{fact}[theorem]{Fact}
\theoremstyle{definition} % non-italics
\iffull
\newtheorem{example}{Example}
\newtheorem{remark}[theorem]{Remark}
\fi
\newtheorem{construction}[theorem]{Construction}

%%%%%%%%%%%%%%%%%%%%%%%%%%%%%%%%%%%%%%%%%%%%%%%%%%%%%%%%%%%%%%%%%%%%%%%%%%%%%%%%



%%%%%%%%%%%%%%%%%%%%%%%%%%%%%%%%%%%%%%%%%%%%%%%%%%%%%%%%%%%%%%%%%%%%%%%%%%%%%%%%
% FORMAT
\newcommand{\keywords}[1]{\bigskip\par\noindent{\footnotesize\textbf{Keywords\/}: #1}}
\newcommand{\FormatAuthor}[3]{
\begin{tabular}{c}
#1 \\ {\small\texttt{#2}} \\ {\small #3}
\end{tabular}
}
\newcommand{\doclearpage}{%
\iffull
\clearpage
\fi
}
\newcommand{\defemph}[1]{\textbf{\emph{#1}}}
\newcommand{\DoQuote}[1]{``#1''}

\newenvironment{boxfig}[2]{\begin{figure}[#1]\fbox{\begin{minipage}{\linewidth}
				\vspace{0.2em}
				\makebox[0.025\linewidth]{}
				\begin{minipage}{0.95\linewidth}
					{{
							#2 }}
				\end{minipage}
				\vspace{0.2em}
	\end{minipage}}}{\end{figure}}


\newcommand{\pprotocol}[4]{
	\begin{boxfig}{H}{
			\begin{center}
				\textbf{#1}
			\end{center}
			#4
			\vspace{0.2em} } \caption{\label{#3} #2}
	\end{boxfig}
}
\newcommand{\protocol}[4]{
	\pprotocol{#1}{#2}{#3}{#4} }
%Example:
%\protocol{Header}{Caption}{label}{the protocol}

%%%%%%%%%%%%%%%%%%%%%%%%%%%%%%%%%%%%%%%%%%%%%%%%%%%%%%%%%%%%%%%%%%%%%%%%%%%%%%%%



%%%%%%%%%%%%%%%%%%%%%%%%%%%%%%%%%%%%%%%%%%%%%%%%%%%%%%%%%%%%%%%%%%%%%%%%%%%%%%%%
% MACROS

\renewcommand{\enspace}{\ensuremath{\thinspace}}
% operators
\newcommand{\arkworks}{\texttt{arkworks}}
\newcommand{\norm}[1]{\ensuremath{\left\lVert#1\right\rVert}}
\newcommand{\abs}[1]{\ensuremath{\Big\vert#1\Big\vert}}
\newcommand{\ip}[1]{\ensuremath{\left\langle #1 \right\rangle}}
\newcommand{\floor}[1]{\ensuremath{\left\lfloor#1\right\rfloor}}
\newcommand{\ceil}[1]{\ensuremath{\left\lceil#1\right\rceil}}
\newcommand{\argmax}{\ensuremath{\mathrm{argmax}}}
\newcommand{\argmin}{\ensuremath{\mathrm{argmin}}}
\newcommand{\xor}{\ensuremath{\oplus}}
\newcommand{\rest}{\mathord{\upharpoonright}}
\newcommand{\sign}{\ensuremath{\mathrm{sgn}}}
\newcommand{\sqrbkt}[1]{\left [#1\right]}
\newcommand{\bkt}[1]{\ensuremath{\left (#1\right)}}
\newcommand{\concat}{\|}

% objects
\newcommand{\F}{\ensuremath{\mathbb{F}}}
\newcommand{\K}{\ensuremath{\mathbb{K}}}
\newcommand{\N}{\ensuremath{\mathbb{N}}}
\newcommand{\Z}{\ensuremath{\mathbb{Z}}}
\newcommand{\R}{\ensuremath{\mathbb{R}}}
\newcommand{\CX}{\ensuremath{\mathbb{C}}}
\newcommand{\E}{\ensuremath{\mathbb{E}}}
\newcommand{\pmset}{\ensuremath{\{-1,1\}}}
\newcommand{\poly}{\ensuremath{\mathrm{poly}}}
%\newcommand{\polylog}{\ensuremath{\mathrm{polylog}}}
%\newcommand{\polyloglog}{\ensuremath{\mathrm{polyloglog}}}
\DeclareMathOperator{\polyloglog}{polyloglog}
\DeclareMathOperator{\polylog}{polylog}
%\newcommand{\loglog}{\ensuremath{\mathrm{loglog}}}
\DeclareMathOperator{\loglog}{loglog}
\DeclareMathOperator{\logloglog}{logloglog}
\newcommand{\diam}{\ensuremath{\mathrm{diam}}}
\newcommand{\const}{\ensuremath{\mathrm{const}}}
\newcommand{\bitstar}{\ensuremath{\{0,1,\star\}}}
\newcommand{\pmstar}{\ensuremath{\{-1,1,\star\}}}
\newcommand{\SharpSAT}{\mathrm{\#SAT}}

% ETHs
\newcommand{\tsat}{\mathrm{3SAT}}
\newcommand{\TimeBound}{T}

\newcommand{\Comp}[1]{\overline{#1}}

\newcommand{\Bits}{\{0,1\}}
\newcommand{\Set}{S}

% oracles
\newcommand{\oracle}{A}
\newcommand{\OracleFamily}{\mathcal{\oracle}}
\newcommand{\OracleDist}[1]{\OracleFamily_{#1}}
%
\newcommand{\uoracle}{F}
\newcommand{\UniSet}{\mathcal{U}}
\newcommand{\negl}{\mathsf{negl}}
% complexity symbols
\newcommand{\NTIME}{\mathrm{NTIME}}
\newcommand{\DTIME}{\mathrm{DTIME}}
\newcommand{\BPTIME}{\mathrm{BPTIME}}
\newcommand{\SPACE}{\mathrm{SPACE}}
\newcommand{\PCP}{\mathrm{PCP}}
\newcommand{\IOP}{\mathrm{IOP}}
\newcommand{\IP}{\mathrm{IP}}
\newcommand{\NP}{\mathrm{NP}}
\newcommand{\PClass}{\mathrm{P}}
\newcommand{\AM}{\mathrm{AM}}
\newcommand{\PSPACE}{\mathrm{PSPACE}}

% complexity classes
\newcommand{\nt}[1]{\NTIME(#1)}
\newcommand{\pcp}[1]{\PCP(#1)}
\newcommand{\dt}[1]{\DTIME(#1)}

% complexity measures
\newcommand{\Time}{\mathsf{t}}
\newcommand{\Event}{\mathsf{E}}
\newcommand{\MachineTime}{T}

\newcommand{\Proof}{\pi}

% malicious
\newcommand{\Malicious}[1]{\tilde{#1}}

% language and relation
\newcommand{\Relation}{\mathcal{R}}
\newcommand{\Language}{L}
%\newcommand{\Index}{\mathbbm{i}}
\newcommand{\Instance}{\mathbbm{x}}
\newcommand{\ExplicitInput}{\mathbbm{x}}
\newcommand{\ImplicitInput}{\mathbbm{y}}
\newcommand{\Witness}{\mathbbm{w}}
\newcommand{\GetLanguage}[1]{\Language(#1)}
\newcommand{\Formula}{\phi}

\newcommand{\Algorithm}{\mathbf{A}}
\newcommand{\Verifier}{\mathbf{V}}
\newcommand{\Prover}{\mathbf{P}}
\newcommand{\Indexer}{\mathbf{I}}
\newcommand{\Decoder}{\mathbf{D}}
\newcommand{\Simulator}{\mathbf{S}}

\newcommand{\MaliciousProof}{\Malicious{\Proof}}
\newcommand{\IndexerProof}{\pi}
\newcommand{\MaliciousIndexerProof}{\Malicious{\IndexerProof}}

\newcommand{\MaliciousProver}{\Malicious{\Prover}}

\newcommand{\DecisionSubscript}{{\scriptscriptstyle \mathsf{dc}}}
\newcommand{\InteractionSubscript}{{\scriptscriptstyle \mathsf{int}}}

\newcommand{\ProverMessage}{m}
\newcommand{\VerifierMessage}{a}

\newcommand{\VerifierQueryAlg}{\Verifier^{\scriptscriptstyle \mathsf{qry}}}
\newcommand{\VerifierDecisionAlg}{\Verifier^{\scriptscriptstyle \mathsf{dc}}}


\newcommand{\VerifierRandomness}{\alpha}
%\newcommand{\DecisionRandomness}[1][]{\Randomness_{\scriptscriptstyle #1 \DecisionSubscript}}
%\newcommand{\InteractionRandomness}[1][]{\Randomness_{\scriptscriptstyle #1 \InteractionSubscript}}
%\newcommand{\DecisionVerifierRandomness}{\VerifierRandomness_{\DecisionSubscript}}
%\newcommand{\DecisionRandomnessLength}{\Randomness_{\Tiny{\mathsf{dc}}}}
%\newcommand{\InteractionRandomnessLength}{\Randomness_{\Tiny{\mathsf{int}}}}

\newcommand{\Tiny}[1]{{\scriptscriptstyle #1}}
\newcommand{\Completeness}{\alpha}
\newcommand{\Length}{\mathsf{l}}
\newcommand{\FuncNum}{\ell}
\newcommand{\SeedLength}{\mathsf{w}}
\newcommand{\LengthTotal}{\mathsf{l}_{\Tiny{\mathsf{tot}}}}
\newcommand{\Queries}{\mathsf{q}}
\newcommand{\RoundsQueried}{\mathsf{q}_{\Tiny{\mathsf{round}}}}
\newcommand{\QueriesMessages}{\Queries_{\FuncNum}}
\newcommand{\Randomness}{\mathsf{r}}
\newcommand{\RandomnessTotal}{\Randomness_{\Tiny{\mathsf{tot}}}}
\newcommand{\Round}{\mathsf{k}}
\newcommand{\Soundness}{\beta}
\newcommand{\Knowledge}{\kappa}
\newcommand{\Alphabet}{\Sigma}
\newcommand{\AlphabetSize}{\lambda}
%\newcommand{\Proximity}{\delta}
\newcommand{\Robustness}{\sigma}
\newcommand{\RBRSoundness}{\beta_{ \Tiny{\mathrm{rbr}}}}
\newcommand{\Advice}{z}
\newcommand{\AdviceSet}{Y}
\newcommand{\AdviceLength}{\ell}
\newcommand{\Error}{\varepsilon}
\newcommand{\ErrorCustom}[1]{\Error_{#1}}
\newcommand{\Tree}{T}
\newcommand{\TreeRoot}{\mathsf{root}}
\newcommand{\TreeValue}{\mathsf{val}}
\newcommand{\EntSample}{z}
\newcommand{\LengthV}{\Length_{\Tiny{\Verifier}}}
\newcommand{\LengthP}{\Length_{\Tiny{\Prover}}}

\newcommand{\TimeI}{\mathsf{it}}
\newcommand{\TimeP}{\mathsf{pt}}
\newcommand{\TimeV}{\mathsf{vt}}
\newcommand{\TimeD}{\mathsf{dt}}
\newcommand{\TimeVDecision}{\mathsf{d}}
\newcommand{\TimeE}{\mathsf{et}}


\newcommand{\Machine}{M}
\newcommand{\MachineSet}{\mathbf{M}}
\newcommand{\MachineRandomness}{r}
\newcommand{\OracleLanguage}{\mathcal{L}}

\newcommand{\cU}{\mathcal{U}}

\newcommand{\SecParam}{\lambda}
\newcommand{\SodParam}{s}

% random oracle
\newcommand{\ROfunc}{\zeta}
\newcommand{\ROquery}{t}
\newcommand{\RODistribution}[1]{\mathcal{U}(#1)}
\newcommand{\GetQueries}[2]{\mathsf{queries}(#1,#2)}
\newcommand{\ROinterface}{\left[\ROfunc\right]}
\newcommand{\ROSpecialSymbol}{\mathsf{set}}


\newcommand{\RepetitionParameter}{t}

% salted soundness
\newcommand{\SaltedSoundnessGame}[4]{\mathsf{SaltedSoundessGame}(#1,#2,#3,#4)}
\newcommand{\QAList}{\sigma}
\newcommand{\SeenOracleStates}{S_{\mathrm{s}}}
\newcommand{\ARGSaltSoundness}{\ARGSoundness_{\mathrm{s}}}

% interactive oracle proof
\newcommand{\IOPSymbol}{\mathsf{IOP}}
\newcommand{\IOPSubscipt}{{\Tiny{\mathsf{IOP}}}}
\newcommand{\IOPIndexer}{\Indexer_{\IOPSubscipt}}
\newcommand{\IOPProver}{\Prover_{\IOPSubscipt}}
\newcommand{\IOPVerifier}{\Verifier_{\IOPSubscipt}}
\newcommand{\IOPDecoder}{\Decoder_{\IOPSubscipt}}
\newcommand{\IOPSystemWI}{(\IOPProver,\IOPVerifier)}
\newcommand{\IOPSystem}{(\IOPProver,\IOPVerifier)}
\newcommand{\IOPVerifierRandomness}{\VerifierRandomness}
\newcommand{\IOPVerifierAllRandomness}{\boldsymbol{\IOPVerifierRandomness}}

\newcommand{\IOPProof}{\Proof}
\newcommand{\Transcript}{\mathsf{tr}}
\newcommand{\IOPTranscript}{\mathsf{tr}}


\newcommand{\IOPQueryDistribution}{\mathcal{D}}
\newcommand{\IOPQueryEntropy}{\mathsf{h}}
\newcommand{\IOPQueryLocations}{\mathcal{I}}
\newcommand{\IOPRandomnessQueryDistribution}{\mathcal{D}}
\newcommand{\IOPRandomnessQueryEntropy}{\mathsf{h}}
%
\newcommand{\IOPAlphabet}{\Alphabet_{\IOPSubscipt}}
\newcommand{\IOPAlphabetSize}{\AlphabetSize_{\IOPSubscipt}}
\newcommand{\IOPQueries}{\Queries_{\IOPSubscipt}}
\newcommand{\IOPRQueries}{\Queries_{\scriptscriptstyle \IOPSymbol, \InteractionSubscript}}
\newcommand{\IOPLength}{\Length_{\IOPSubscipt}}
\newcommand{\IOPRandomness}{\Randomness_{\IOPSubscipt}}

\newcommand{\IOPCompleteness}{\Completeness_{\IOPSubscipt}}
\newcommand{\IOPSoundness}{\Soundness_{\IOPSubscipt}}
\newcommand{\IOPRBRSoundness}{\Soundness_{\scriptscriptstyle \IOPSymbol,\mathrm{rbr}}}
\newcommand{\IOPKnowledge}{\Knowledge_{\IOPSubscipt}}
\newcommand{\IOPIndexerRandomness}{\rho_{0}}
%
\newcommand{\IOPTimeI}{\TimeI_{\IOPSubscipt}}
\newcommand{\IOPTimeP}{\TimeP_{\IOPSubscipt}}
\newcommand{\IOPTimeV}{\TimeV_{\IOPSubscipt}}
%
\newcommand{\IOPInteractionRandomness}{\InteractionRandomness[\IOPSubscipt ,]}
\newcommand{\IOPDecisionRandomness}{\DecisionRandomness[ \IOPSubscipt ,]}
\newcommand{\IOPVerifierDecisionRandomness}{\DecisionVerifierRandomness}
\newcommand{\IOPMaliciousProver}{\MaliciousProver_{\IOPSubscipt}}
\newcommand{\IOPMaliciousIndexerProof}{\MaliciousIndexerProof}
\newcommand{\IOPMaliciousProof}{\MaliciousProof}

% IP protocol
\newcommand{\IPSymbol}{\mathsf{IP}}
\newcommand{\IPSubscipt}{{\Tiny{\IPSymbol}}}
\newcommand{\IPProver}{\Prover_{\IPSubscipt}}
\newcommand{\IPVerifier}{\Verifier_{\IPSubscipt}}
\newcommand{\IPProverToVerifier}{\Prover \to \Verifier}
\newcommand{\IPMaliciousProver}{\MaliciousProver_{\IPSubscipt}}
\newcommand{\IPCompleteness}{\Completeness_{\IPSubscipt}}
\newcommand{\IPSoundness}{\Soundness_{\IPSubscipt}}
\newcommand{\IPLength}{\Length_{\IPSubscipt}}
\newcommand{\IPRandomness}{\Randomness_{\IPSubscipt}}
\newcommand{\IPRound}{\Round_{\IPSubscipt}}
\newcommand{\IPVerifierRandomness}{\VerifierRandomness}
\newcommand{\IPTimeV}{\TimeV_{\IPSubscipt}}
\newcommand{\IPDecisionRandomness}{\DecisionRandomness[ \IPSymbol ,]}
\newcommand{\IPVerifierDecisionRandomness}{\DecisionVerifierRandomness}
\newcommand{\IPRQueries}{\Queries_{\scriptscriptstyle \IPSymbol, \InteractionSubscript}}
\newcommand{\IPInteractionRandomness}{\InteractionRandomness[\IPSymbol ,]}

% pIOP protocol
\newcommand{\FuncNumCompiler}{s}
\newcommand{\PIOP}{\textrm{poly}-IOP\xspace}
\newcommand{\PIOPP}{\textrm{poly}-IOPP\xspace}
\newcommand{\PIOPSymbol}{\mathsf{poly}}
\newcommand{\PIOPSubscript}{{\Tiny{\PIOPSymbol}}}
\newcommand{\PIOPProver}{\Prover_{\PIOPSubscript}}
\newcommand{\PIOPVerifier}{\Verifier_{\PIOPSubscript}}
\newcommand{\PIOPMaliciousProver}{\MaliciousProver_{\PIOPSubscript}}
\newcommand{\PIOPCompleteness}{\Completeness_{\PIOPSubscript}}
\newcommand{\PIOPSoundness}{\Soundness_{\PIOPSubscript}}
\newcommand{\PIOPLength}{\Length_{\PIOPSubscript}}
\newcommand{\PIOPFuncNum}{\FuncNumCompiler_{\PIOPSubscript}}
\newcommand{\PIOPLogFuncNum}{{\mu_{\PIOPSubscript}}}
\newcommand{\PIOPRandomness}{\Randomness_{\PIOPSubscript}}
\newcommand{\PIOPRound}{\Round_{\PIOPSubscript}}
\newcommand{\PIOPTimeV}{\TimeV_{\PIOPSubscript}}
\newcommand{\PIOPTimeE}{\TimeE_{\PIOPSubscript}}
\newcommand{\PIOPQueriesInput}{\Queries_{\PIOPSubscript, \Tiny{\ImplicitInput}}}
\newcommand{\PIOPQueriesProof}{\Queries_{\PIOPSubscript, \Tiny{\Proof}}}
\newcommand{\PIOPQueriesMessages}{\Queries_{\PIOPSubscript, \Tiny{\FuncNum}}}
\newcommand{\PIOPState}{\State_{\PIOPSubscript}}
\newcommand{\PIOPVerifierRandomness}[1]{\VerifierRandomness_{#1}}

\newcommand{\IOPRound}{\PIOPRound}

% PCP
\newcommand{\PCPSymbol}{\mathsf{PCP}}
\newcommand{\PCPSubscipt}{{\Tiny{\PCPSymbol}}}
\newcommand{\PCPProver}{\Prover_{\PCPSubscipt}}
\newcommand{\PCPVerifier}{\Verifier_{\PCPSubscipt}}
\newcommand{\PCPSystem}{(\PCPProver,\PCPVerifier)}

% hashing
\newcommand{\Universe}{U}
\newcommand{\HashFunction}{h}
\newcommand{\HashFamily}{\mathcal{H}}
\newcommand{\HashTime}{\mathsf{ht}}
\newcommand{\Bin}{b}

% compiler construction
\newcommand{\TableT}{T}
\newcommand{\TableSize}{M}
\newcommand{\Xqueries}{X}

\usepackage{soul}
\newcommand{\highlight}[1]{\hl{#1}}

% transformation
\newcommand{\Transformation}{\mathbb{T}}

%entropy
\newcommand{\supp}{\mathsf{supp}}
\newcommand{\ShanEnt}{H}
\newcommand{\MinEnt}{\ShanEnt_{\mathsf{min}}}
\newcommand{\EE}[2]{\mathop{\mathbf E}_{#1}\! \left[ {#2} \right]}

\newcommand{\PRG}{\mathsf{G}}
\newcommand{\PRGSubscript}{\Tiny{\mathsf{PRG}}}
\newcommand{\PRGSeedLength}{\ell_{\PRGSubscript}}
\newcommand{\PRGCircuitSize}{\mathsf{s}_{\PRGSubscript}}
\newcommand{\PRGError}{\epsilon_{\PRGSubscript}}
\newcommand{\PRGRunningTime}{\mathsf{t}_{\PRGSubscript}}
\newcommand{\PRGTime}{\mathsf{t}_{\PRGSubscript}}

\newcommand{\CSPInstance}{\psi}
\newcommand{\CSPAlphabet}{\Alphabet}
\newcommand{\CSPArity}{\Queries}
\newcommand{\CSPLength}{\Length}
\newcommand{\CSPVariables}{n}
\newcommand{\CSPConstraintNum}{m}
\newcommand{\CSPConstraint}{C}
\newcommand{\CSPConstraintFunction}{f}
\newcommand{\CSPAssignment}{a}

\newcommand{\class}[1]{\mathsf{#1}}
\newcommand{\EClass}{{\class{E}}}
\newcommand{\Interaction}[2]{\langle #1, #2 \rangle}
\newcommand{\Circuit}{C}
\newcommand{\CircuitPoly}{\LDE{C}}
\newcommand{\CircuitFamily}{\mathsf{Circ}}
\newcommand{\AmplificationParam}{t}

%Macros Amey
\newcommand{\calK}{\mathcal{K}}

% table coloring
\definecolor{paleyellow}{RGB}{255, 255, 167}
\newcommand{\HighlightRow}{\rowcolor{paleyellow}}

% length and round reduciton
\newcommand{\SampleSet}{I}
\newcommand{\SampleSize}{\ell}

% perfect completeness
\newcommand{\PerfectCompAdviceSize}{t}

% limits for public coin iops
\newcommand{\LenLimitFunction}{m}

% limits for private coin ips
\newcommand{\Rej}{\mathsf{Rej}}
\newcommand{\IPLimClusterNum}{t}
\newcommand{\IPLimCluster}{C}
\newcommand{\IPLimClusterDiff}{\varepsilon}
\newcommand{\IPLimSetSize}{N}
\newcommand{\IPLimRepetitions}{w}
\newcommand{\SLBSymbol}{\Tiny{\mathsf{SLB}}}
\newcommand{\SLBProver}{\Prover_{\SLBSymbol}}
\newcommand{\SLBMaliciousProver}{\MaliciousProver_{\SLBSymbol}}
\newcommand{\SLBVerifier}{\Verifier_{\SLBSymbol}}

% ETH assumptions
\newcommand{\ETH}{\mathtt{ETH}}
\newcommand{\RETH}{\mathtt{RETH}}
\newcommand{\AMETH}{\mathtt{AMETH}}


\newcommand{\Polynomial}{p}
\newcommand{\Deg}{\mathsf{deg}}
\newcommand{\Degree}{d}
\newcommand{\Domain}{D}
\newcommand{\Distance}{\delta}
\newcommand{\Rate}{\rho}
\newcommand{\ReedSolomon}{\mathsf{RS}}
\newcommand{\ConstrainedReedSolomon}{\mathsf{CRS}}
\newcommand{\InterleavedReedSolomon}{\mathsf{IRS}}
\newcommand{\ReedMuller}{\mathsf{RM}}
\newcommand{\ListDecoder}{\mathsf{ListDec}}
\newcommand{\Generator}{g}
\newcommand{\GeneratedSet}[1]{\langle #1 \rangle}
\newcommand{\SummationDomain}{H}
\newcommand{\EvaluationDomain}{\mathcal{L}}
\newcommand{\VanishingPolynomial}{\LDE{V}}
\newcommand{\LDE}[1]{\hat{#1}}
\newcommand{\LDEWide}[1]{\widehat{#1}}
%\newcommand{\AgreementDomain}{\mathsf{AgrDom}}
\newcommand{\Agreement}{\mu}
\newcommand{\ListDecode}{\mathsf{ListDec}}
\newcommand{\CorrelatedAgreement}{\mathsf{CorrAgr}}
\newcommand{\CorrelatedAgreementSize}{\mathsf{CorrAgrSize}}
\newcommand{\Code}{\mathcal{C}}
\newcommand{\CodeWithDeg}[2]{\Code^{\Tiny{(#2)}}_{{#1}}}

\newcommand{\QueriesImplicitInput}{\Queries_{\Tiny{\ImplicitInput}}}

\newcommand{\QueriesCode}{\Queries_{\Tiny{\Code}}}
\newcommand{\QueriesString}{\Queries_{\Tiny{\mathsf{s}}}}
\newcommand{\RoundCode}{\Round_{\Tiny{\Code}}}

%%%%% List Decoding $$$$$
\newcommand{\List}{\Lambda}
\newcommand{\ListSize}{\ell}
\newcommand{\ListSizeInterleaved}[1]{\ell_{#1}}
\newcommand{\ListProximity}{\gamma}

%%%%% Low Degree Test %%%%%
\newcommand{\LDTSymbol}{\Tiny{\mathsf{prx}}}
\newcommand{\LDTQueries}{\Queries_{\LDTSymbol}}
\newcommand{\LDTQueriesInput}{\Queries_{\LDTSymbol,f}}
\newcommand{\LDTQueriesProof}{\Queries_{\LDTSymbol,\Proof}}
\newcommand{\LDTRounds}{\Round_{\LDTSymbol}}
\newcommand{\LDTRandomness}{\Randomness_{\LDTSymbol}}
\newcommand{\LDTSoundness}{\Soundness_{\LDTSymbol}}
\newcommand{\LDTVerifier}{\Verifier_{\LDTSymbol}}
\newcommand{\LDTProver}{\Prover_{\LDTSymbol}}
\newcommand{\LDTLength}{\Length_{\LDTSymbol}}
\newcommand{\LDTCode}{\Code_{\LDTSymbol}}
\newcommand{\LDTRate}{\Rate_{\LDTSymbol}}
\newcommand{\LDTTimeV}{\TimeV_{\LDTSymbol}}
\newcommand{\LDTAlphabet}{\Alphabet_{\LDTSymbol}}
\newcommand{\QueriesFunction}{\Queries_{\Tiny{f}}}
\newcommand{\QueriesProof}{\Queries_{\Tiny{\Proof}}}
\newcommand{\LDTDegree}{\Degree_{\LDTSymbol}}
\newcommand{\LDTState}{\State_{\LDTSymbol}}
\newcommand{\LDTVerifierRandomness}[1]{\VerifierRandomness_{\Tiny{\LDTSymbol, #1}}}
\newcommand{\LDTTranscript}{\Transcript_{\LDTSymbol}}


%%%%% Constrained RS Test %%%%%
\newcommand{\CRSSymbol}{\Tiny{\mathsf{prx}}}
\newcommand{\CRSQueries}{\Queries_{\CRSSymbol}}
\newcommand{\CRSQueriesInput}{\Queries_{\CRSSymbol,f}}
\newcommand{\CRSQueriesProof}{\Queries_{\CRSSymbol,\Proof}}
\newcommand{\CRSRounds}{\Round_{\CRSSymbol}}
\newcommand{\CRSRandomness}{\Randomness_{\CRSSymbol}}
\newcommand{\CRSSoundness}{\Soundness_{\CRSSymbol}}
\newcommand{\CRSVerifier}{\Verifier_{\CRSSymbol}}
\newcommand{\CRSProver}{\Prover_{\CRSSymbol}}
\newcommand{\CRSLength}{\Length_{\CRSSymbol}}
\newcommand{\CRSCode}{\Code_{\CRSSymbol}}
\newcommand{\CRSRate}{\Rate_{\CRSSymbol}}
\newcommand{\CRSTimeV}{\TimeV_{\CRSSymbol}}
\newcommand{\CRSAlphabet}{\Alphabet_{\CRSSymbol}}
\newcommand{\CRSDegree}{\Degree_{\CRSSymbol}}
\newcommand{\CRSState}{\State_{\CRSSymbol}}
\newcommand{\CRSVerifierRandomness}[1]{\VerifierRandomness_{\Tiny{\CRSSymbol, #1}}}
\newcommand{\CRSTranscript}{\Transcript_{\CRSSymbol}}



%%%%%% Polynomials %%%%%%
\newcommand{\EvalPoint}{a}
\newcommand{\EvalAns}{b}
\newcommand{\EvalPointAlt}{x}
\newcommand{\EvalAnsAlt}{y}
\newcommand{\EvalHoleFill}{c}
\newcommand{\Coefficient}{r}
\newcommand{\CoefficientVector}{\vec{\Coefficient}}
\newcommand{\CoefficientVectorSubscript}{\Tiny{\vec{\Coefficient}}}

\newcommand{\Quotient}{\mathsf{Quotient}}
\newcommand{\PolyQuotient}{\mathsf{PolyQuotient}}

\newcommand{\Unquotient}{\mathsf{Unquotient}}




\newcommand{\IndicatorPolynomial}{I}

\newcommand{\ProximitySubscript}{\Tiny{\LDTSymbol}}
\newcommand{\ProximityTestDomain}{\EvaluationDomain_{\LDTSymbol}}

\newcommand{\OutOfDomainSymbol}{\mathsf{out}}
\newcommand{\OutOfDomainDomain}{L_{\Tiny{\OutOfDomainSymbol}}}
\newcommand{\OutOfDomainPoint}{x}
\newcommand{\OutOfDomainAns}{y}
\newcommand{\OutOfDomainHoleFill}{z}


\newcommand{\BivariateSymbol}{\mathsf{Bivar}}
\newcommand{\BivariateVerifier}{\Verifier_{\Tiny{\BivariateSymbol}}}
\newcommand{\BivariateProver}{\Prover_{\Tiny{\BivariateSymbol}}}
\newcommand{\BivariateError}{\Error_{\Tiny{\BivariateSymbol}}}


\newcommand{\DegreeMax}{d_{\Tiny{\mathsf{max}}}}
\newcommand{\DegreeMin}{d_{\Tiny{\mathsf{min}}}}
\newcommand{\RateCombined}{\Rate_{\Tiny{X+Y}}}
\newcommand{\RateMax}{\Rate_{\Tiny{\mathsf{max}}}}
\newcommand{\EvaluationDomainMatrix}{\Domain_{f}}
\newcommand{\CodeMatrix}{\Code_{f}}
\newcommand{\CodeColumn}{\Code_{\Tiny{X}}}
\newcommand{\DegreeRow}{\Degree_{\Tiny{Y}}}
\newcommand{\DegreeColumn}{\Degree_{\Tiny{X}}}
\newcommand{\RateColumn}{\Rate_{\Tiny{X}}}
\newcommand{\EvaluationDomainRow}{\EvaluationDomain_{\Tiny{Y}}}
\newcommand{\EvaluationDomainColumn}{\EvaluationDomain_{\Tiny{X}}}
\newcommand{\EvaluationDomainExtension}{\EvaluationDomain_{\mathsf{ext}}}
\newcommand{\ProximityGapsSubscript}{\Tiny{\star}}
\newcommand{\ProximityGenerator}{\mathsf{Gen}}
\newcommand{\ProximityGeneratorError}{\ErrorCustom{\ProximityGapsSubscript}}
\newcommand{\ProximityGeneratorBound}{\ProximityBound_{\ProximityGapsSubscript}}
\newcommand{\ProximityGeneratorSeed}{\alpha}
\newcommand{\ProximityGeneratorSeedLength}{\SeedLength_{\ProximityGapsSubscript}}
\newcommand{\ProximityGeneratorTime}{\Time_{\ProximityGapsSubscript}}
\newcommand{\ProximityGapsFuncNum}{\FuncNum_{\ProximityGapsSubscript}}
\newcommand{\ROCS}{\mathsf{R1CS}}
\newcommand{\ROCSRelation}{\Relation_{\Tiny{\ROCS}}}
\newcommand{\GROCS}{\mathsf{GR1CS}}
\newcommand{\GROCSRelation}{\Relation_{\Tiny{\GROCS}}}
\newcommand{\SummationDomainInput}{\SummationDomain_{\Tiny{\mathsf{in}}}}
\newcommand{\SummationDomainAux}{\SummationDomain_{\Tiny{\mathsf{aux}}}}
\newcommand{\PowerPolynomial}{\LDE{p}}
\newcommand{\MatrixPolynomial}{\LDE{q}}

\newcommand{\Order}{\mathsf{Ord}}
\newcommand{\QuotientAns}{\mathsf{Ans}}
\newcommand{\QuotientFill}{\mathsf{Fill}}

\newcommand{\EventOutOfDomain}{\Event_{\Tiny{\OutOfDomainSymbol}}}
\newcommand{\EventLDT}{\Event_{\Tiny{\LDTSymbol}}}
\newcommand{\Query}{t}
\newcommand{\DegreeIntro}{\Degree^{*}}

\newcommand{\State}{\mathsf{State}}
\newcommand{\ProofOfWorkLength}{w}
\newcommand{\ProofOfWorkMessage}{\VerifierRandomness_{\Tiny{\mathsf{w}}}}


\newcommand{\Iterations}{M}

\newcommand{\SetS}{\mathcal{S}}
\newcommand{\SetW}{\mathcal{W}}
\newcommand{\SetG}{\mathcal{G}}
\newcommand{\SetD}{\mathcal{D}}
\newcommand{\SetL}{\mathcal{L}}
\newcommand{\SetZ}{\mathcal{Z}}
\newcommand{\SetP}{\mathcal{P}}
\newcommand{\PolynomialVariable}{X}
\newcommand{\PolynomialVariableTwo}{Y}
\newcommand{\Variable}{x}
\newcommand{\DegreeInitial}{\Degree_{0}}
\newcommand{\EvaluationDomainInitial}{\EvaluationDomain_{0}}
\newcommand{\EvaluationDomainSize}{n}

\newcommand{\Combine}{\mathsf{Combine}}

\newcommand{\Func}{f}
\newcommand{\DegreeCorrection}{e}
\newcommand{\DegreeLower}{e}
\newcommand{\FuncBivariate}{Q}
\newcommand{\Index}{j}
\newcommand{\IterationIndex}{i}

\newcommand{\Row}{r}

\newcommand{\IntersectionPointSet}{\mathcal{B}}


\newcommand{\ErrorListDecoding}{\Error^{\Tiny{\mathsf{list}}}}
\newcommand{\ErrorProximityGaps}{\mathsf{err}^{\ProximityGapsSubscript}}
\newcommand{\ErrorProximityGapsMany}{\mathsf{err}}
\newcommand{\ProximityGapsBound}{\mathsf{B}^{\ProximityGapsSubscript}}
\newcommand{\ErrorProximityGapsWeak}{\mathsf{err}}
\newcommand{\ProximityGapsBoundWeak}{\mathsf{B}}


\newcommand{\Codeword}{c}
\newcommand{\CodeRow}{\Code^{g}}
\newcommand{\CodeCoefficients}{\Code^{\Tiny{\mathsf{coef}}}}


\newcommand{\TranscriptPrevious}{\Transcript_{\Tiny{\mathsf{prev}}}}
\newcommand{\TranscriptCurrent}{\Transcript_{\Tiny{\mathsf{cur}}}}

\newcommand{\SplitToBivariate}{\mathsf{Split}}

\newcommand{\Fill}{\mathsf{Fill}}
\newcommand{\Ans}{\mathsf{Ans}}
\newcommand{\DegreeCorrect}{\mathsf{DegCor}}

% Shake specific things
\newcommand{\NumeratorPoly}{\LDE{\mathsf{u}}}
\newcommand{\DenominatorPoly}{\LDE{\mathsf{w}}}
\newcommand{\Zeros}{\mathsf{zeros}}
\newcommand{\ErrorRandomness}{\Error^{\Tiny{\mathsf{rand}}}}
\newcommand{\ErrorProximity}{\Error^{\Tiny{\mathsf{prox}}}}
\newcommand{\QueryFunction}{\phi}
\newcommand{\QuerySet}{\mathcal{Q}}
\newcommand{\Array}{A}


\newcommand{\FoldingParameter}{k}
\newcommand{\FoldSymbol}{\Tiny{\mathsf{fold}}}
\newcommand{\STIRSymbol}{\Tiny{\mathsf{shift}}}
\newcommand{\OutSymbol}{\Tiny{\mathsf{out}}}
\newcommand{\FinalSymbol}{\Tiny{\mathsf{fin}}}
\newcommand{\EvaluationDomainFolded}{\mathcal{D}}
\newcommand{\Point}{z}
\newcommand{\PointAns}{y}
\newcommand{\RandomnessSTIR}{\Point}
\newcommand{\RandomnessOut}{r^{\OutSymbol}}
\newcommand{\RandomnessFold}{\alpha}

\newcommand{\OODRepetition}{s}

\newcommand{\CombineSymbol}{\Tiny{\mathsf{comb}}}
\newcommand{\RandomnessCombine}{\gamma}

\newcommand{\Fold}{\mathsf{Fold}}
%\newcommand{\FoldPoly}{\mathsf{PolyFold}}

\newcommand{\RandomnessFinal}{r^{\FinalSymbol}}
\newcommand{\RandomnessFinalBold}{\vect{r}^{\FinalSymbol}}
\newcommand{\OutAns}{\beta}
\newcommand{\ErrorFoldInitial}{\Error^{\FoldSymbol}}
\newcommand{\ErrorOut}{\Error^{\OutSymbol}}
\newcommand{\ErrorSTIR}{\Error^{\STIRSymbol}}
\newcommand{\ErrorFold}{\Error^{\FoldSymbol}}
\newcommand{\ErrorFinal}{\Error^{\Tiny{\mathsf{fin}}}}
\newcommand{\FuncShifted}{g}
\newcommand{\StoppingDegree}{\Degree_{\Tiny{\mathsf{stop}}}}

\newcommand{\ErrorROCSShift}{\Error^{\Tiny{\mathsf{shift}}}}
\newcommand{\ErrorROCSDec}{\Error^{\Tiny{\mathsf{dec}}}}

\newcommand{\ErrorCompilerOut}{\Error^{\Tiny{\mathsf{out}}}}
\newcommand{\ErrorCompilerPIOP}{\Error^{\Tiny{\mathsf{piop}}}}
\newcommand{\ErrorCompilerCombine}{\Error^{\Tiny{\mathsf{com}}}}
\newcommand{\ErrorCompilerLDT}{\Error^{\Tiny{\mathsf{prx}}}}
\newcommand{\ErrorLDT}[1]{\mathsf{err}^{\Tiny{\LDTSymbol}}_{#1}}
\newcommand{\ErrorPIOP}[1]{\mathsf{err}^{\PIOPSubscript}_{#1}}
\newcommand{\ErrorCompilerSumcheck}{\Error^{\Tiny{\mathsf{sum}}}}

\newcommand{\Extractor}{\mathbf{E}}
\newcommand{\PIOPExtractor}{\Extractor_{\PIOPSubscript}}

\newcommand{\ExtractoRS}{\Extractor_{\Tiny{\mathsf{RS}}}}
\newcommand{\ExtractoRSTime}{\TimeE_{\Tiny{\mathsf{RS}}}}
\newcommand{\DegreeFix}{\Degree^{*}}
\newcommand{\ExtractoCRS}{\Extractor_{\Tiny{\mathsf{CRS}}}}
\newcommand{\ExtractoCRSTime}{\TimeE_{\Tiny{\mathsf{CRS}}}}

\newcommand{\PCS}{\mathsf{PCS}}
\newcommand{\PCSCommit}{\PCS.\mathsf{Commit}}
\newcommand{\PCSOpen}{\PCS.\mathsf{Open}}
\newcommand{\PCSEvaluate}{\PCS.\mathsf{Eval}}
\newcommand{\PCSVerify}{\PCS.\mathsf{Ver}}
\newcommand{\PCSEvaluateVerify}{\PCS.\mathsf{EvalVer}}
\newcommand{\NumVariables}{m}
\newcommand{\NumVariablesInput}{{m_{\scriptscriptstyle \mathrm{in}}}}
\newcommand{\NumVariablesWitness}{{m_{\scriptscriptstyle \mathrm{wit}}}}
\newcommand{\EQPoly}{\mathsf{eq}}
\newcommand{\MV}[1]{\mathsf{mv}[{#1}]}
\newcommand{\UV}[1]{\mathsf{uv}[{#1}]}
\newcommand{\coeffUV}{\mathsf{coeff}_{\scriptscriptstyle \mathsf{UV}}}
\newcommand{\coeffMV}{\mathsf{coeff}_{\scriptscriptstyle \mathsf{MV}}}
\newcommand{\vect}[1]{{\bm{#1}}}
\newcommand{\PowMap}{\mathsf{pow}}
\newcommand{\coeffVec}{\vect{a}}

\newcommand{\OODPoint}{\Point}
\newcommand{\OODAns}{y}
\newcommand{\NumPoints}{{\RepetitionParameter_{0}}}
\newcommand{\SumcheckPoly}{\LDE{h}}
\newcommand{\EQCombinePoly}{\LDE{w}}
\newcommand{\InitPoint}{x}
\newcommand{\InitAns}{y}
\newcommand{\Target}{\sigma}
\newcommand{\Block}{\mathsf{Block}}
\newcommand{\Coset}{\Block}
\newcommand{\BlockDistance}{\Delta_{\Tiny{\mathsf{b}}}}
\newcommand{\CosetDistance}{\BlockDistance}
\newcommand{\BlockList}{\List_{\Tiny{\mathsf{b}}}}
\newcommand{\CosetList}{\BlockList}

\newcommand{\CodeRS}{\Code_{\Tiny{\mathsf{RS}}}}
\newcommand{\CodeRSFold}{\Code_{\Tiny{\mathsf{FRS}}}}
\newcommand{\CodeCRS}{\Code_{\Tiny{\mathsf{CRS}}}}
\newcommand{\CodeCRSFold}{\Code_{\Tiny{\mathsf{FCRS}}}}
\newcommand{\FoldedReedSolomon}{\mathsf{FRS}}
\newcommand{\FoldedConstrainedReedSolomon}{\mathsf{FCRS}}
\newcommand{\WeightPolynomial}{\LDE{w}}

\newcommand{\QueryWHIR}{q_{\Tiny{\mathsf{WHIR}}}}
\newcommand{\QueryFRI}{q_{\Tiny{\mathsf{FRI}}}}
\newcommand{\QueryBaseFold}{q_{\Tiny{\mathsf{BF}}}}
\newcommand{\QuerySTIR}{q_{\Tiny{\mathsf{STIR}}}}
\newcommand{\QueryX}{q_{\Tiny{\mathsf{X}}}}

\newcommand{\plonk}{\ensuremath{\mathcal{P} \mathfrak{lon}\mathcal{K}}\xspace}
\newcommand{\SumcheckIOPSymbol}{\text{-IOP}}
\newcommand{\SumcheckIOPsSymbol}{\text{-IOPs}}
\newcommand{\SumcheckIOP}{\ensuremath{\Sigma\SumcheckIOPSymbol}\xspace}
\newcommand{\SumcheckIOPTitle}{\texorpdfstring{\SumcheckIOP}{Sigma-IOP}}
\newcommand{\SumcheckIOPs}{\ensuremath{\Sigma\SumcheckIOPsSymbol}\xspace}
\newcommand{\SumcheckIOPP}{\ensuremath{\Sigma\SumcheckIOPSymbol}\xspace}
\newcommand{\SumcheckIOPPTitle}{\texorpdfstring{\SumcheckIOPPTitle}{Sigma-IOPP}}
\newcommand{\SumcheckIOPPs}{\ensuremath{\Sigma\SumcheckIOPsSymbol}\xspace}
\newcommand{\SumcheckIOPPsTitle}{\texorpdfstring{\SumcheckIOPPs}{Sigma-IOPPs}}
\newcommand{\FunctionClassIOPSymbol}{\text{-IOP}}
\newcommand{\FunctionClassIOPsSymbol}{\text{-IOPs}}
\newcommand{\FunctionClassIOPP}{\ensuremath{\mathcal{F}\FunctionClassIOPSymbol}\xspace}
\newcommand{\FunctionClassIOPPs}{\ensuremath{\mathcal{F}\FunctionClassIOPsSymbol}\xspace}
\newcommand{\FunctionClassIOPPsTitle}{\texorpdfstring{\FunctionClassIOPPs}{F-IOPPs}}


\newcommand{\Bijection}{\phi}

%color remapping
\definecolor{graphgreen}{RGB}{44, 160, 44}
\colorlet{green}{graphgreen}
\definecolor{graphblue}{RGB}{31,119, 180}
\colorlet{blue}{graphblue}
\definecolor{graphpurple}{RGB}{148,103,189}
\colorlet{purple}{graphpurple}
\definecolor{graphred}{RGB}{214,39,40}
\colorlet{red}{graphred}


%%%%%%%%%%%%%%%%%%%%%%%%%%%%%%%%%%%%%%%%%%%%%%%%%%%%%%%%%%%%%%%%%%%%%%%%%%%%%%%%
%%%%%%%%%%%%%%%%%%%%%%%%%%%%%%%%%%%%%%%%%%%%%%%%%%%%%%%%%%%%%%%%%%%%%%%%%%%%%%%%
%%%%%%%%%%%%%%%%%%%%%%%%%%%%%%%%%%%%%%%%%%%%%%%%%%%%%%%%%%%%%%%%%%%%%%%%%%%%%%%%
% Commands retained from declarations inside the original template body.
\newcommand{\CircuitDegree}{c}
\newcommand{\SecParamProtocol}{\SecParam_{\scriptscriptstyle\mathsf{prot}}}
\newcommand{\IOPSpec}{\mathsf{Spec}}
\newcommand{\IOPSpecFull}{(\Alphabet, \IOPRound, (\FuncNumCompiler_i)_{i \in [\IOPRound]}, ((\OracleClass_{i, j}, \QueryClass_{i, j}, \QueryFunc_{i, j})_{j \in [\FuncNumCompiler_i]})_{i \in [\IOPRound]})}
\newcommand{\OracleClass}{\mathcal{O}}
\newcommand{\QueryClass}{\mathcal{W}}
\newcommand{\QueryFunc}{\mathcal{F}}
\newcommand{\ConstraintSystem}{\mathcal{C}}

% Notation for this note; the template's mathematical commands are retained.
\newcommand{\Tate}{\mathsf{T}_{8}}
\newcommand{\MTLeaf}{L_{8}}
\newcommand{\MTNode}{C}
\newcommand{\MTCommit}{\mathsf{Commit}}
\newcommand{\MTOpen}{\mathsf{Open}}
\newcommand{\MTVerify}{\mathsf{Verify}}
\newcommand{\MDLeaf}{L_{\mathrm{MD}}}
\newcommand{\mtbits}{\{0,1\}^{n}}
\newcommand{\mtadv}{\operatorname{Adv}}
\renewcommand{\arraystretch}{1.5}
\hypersetup{pdftitle={Faster Merkle Commitments with T8 Leaves},pdfauthor={Thomas Coratger, Giacomo Fenzi, and Eylon Yogev}}

\begin{document}
\title{Faster Merkle Commitments with \texorpdfstring{$\Tate$}{T8} Leaves}
\author{Thomas Coratger \and Giacomo Fenzi \and Eylon Yogev}
\date{}
\maketitle

\begin{abstract}
We describe a simple optimization for Merkle trees whose leaves are records
of eight hash-sized blocks. Hash each record with $\Tate$, the eight-input
extension of T5, and compute the upper tree in the ordinary way. Compared
with a four-call fixed-length leaf hash, commitment to $N$ records uses
$4N-1$ rather than $5N-1$ compression calls, and verification uses
$\log_2N+3$ rather than $\log_2N+4$. The opened record and its authentication
path have exactly the same sizes. We give the construction, a binding
reduction, and precise cost comparisons, including batched openings and
applications to Merkle-based SNARKs. We also describe concrete SHA-256,
BLAKE3, and Keccak experiments: using more of the native compression input,
combining gadgets with short MD tails, and exposing independent work to SIMD.
Measured improvements depend on record width and on the optimized baseline;
the concrete research modes require assumptions beyond compressor collision
resistance alone.
\end{abstract}

\section{The idea}
A Merkle tree over records performs two kinds of hashing: it first compresses
each record to one digest, then combines digests up to the root. An opening
repeats the first computation on the disclosed record and checks the usual
path of sibling digests. Consequently, a more efficient leaf hash can reduce
both commitment and verification work while leaving the upper tree and the
opening format unchanged.

We use the eight-input extension of T5 described by Dodis, Khovratovich,
Mouha, and Nandi in Section~6 of~\cite{DKMN21}. This function, denoted
$\Tate$, hashes eight $n$-bit blocks to one $n$-bit digest using three
$3n$-to-$n$ compression calls. A conventional fixed-length
Merkle--Damg\aa rd leaf hash using the same primitive width needs four calls.
Replacing the latter by $\Tate$ therefore saves one call at each leaf.

The opening unit is central: a position contains an entire eight-block
record, and the verifier receives all eight blocks. Both schemes commit to
the same records. We do not compare this interface with opening one word
in a tree containing eight times as many leaves. Since the verifier knows
every input to $\Tate$, security uses its ordinary collision resistance;
no specialized partial-opening property is needed.

The basic construction concerns the placement and use of an existing
compression construction. The change is confined to the leaves and applies
when complete records are already the unit of access, including suitable
oracle-table commitments used in succinct arguments. The first part of this
note gives an abstract call-count comparison. Section~\ref{sec:concrete}
then records concrete adapters and implementation lessons, including new
research variants whose security status is stated separately.
Section~\ref{sec:measurements} reports measured running times and identifies
the particular standard-hash comparator used in each experiment.

\section{Construction and openings}
\label{sec:construction}
\subsection{Records and hash interfaces}
Fix a digest length $n$ and a public record count $N=2^d$. The committed
vector is
\[
 X=(X_0,\ldots,X_{N-1}),\qquad
 X_i=(x_{i,1},\ldots,x_{i,8})\in(\{0,1\}^n)^8.
\]
Record length, block order, $N$, and the hash-suite identifier belong to the
fixed public context. If an application supports several formats, it must
bind the context as well. Any extra context-binding cost should be included
identically in both schemes.

Let $h_1,h_2,h_3,h_{\mathrm N},h_{\mathrm B}$ be domain-separated
$3n$-to-$n$ compression functions. In the idealized security analysis they
are independent random functions. In an implementation their roles must be
separated using a specified primitive interface. The names do not justify
using one untagged function for all roles.

For a record $Z=(z_1,\ldots,z_8)$, set
\begin{equation}
\begin{aligned}
 a&=h_1(z_1,z_2,z_3),\qquad b=h_2(z_4,z_5,z_6),\\
 L_8(Z)=\Tate(Z)&=h_3(a\oplus z_7,b\oplus z_7,z_8)\oplus z_7.
\end{aligned}
\label{eq:t8}
\end{equation}
This is the larger-input extension from~\cite[Section~6]{DKMN21}, with an
additional input block in each of the three compression calls. There are
three calls and three $n$-bit XORs. The first two calls are independent, so
the compression depth is two when parallel evaluation is available.

For the upper tree, use the ordered binary node hash
\begin{equation}
 C(a,b)=h_{\mathrm N}(a,b,0^n).
 \label{eq:node}
\end{equation}
The fixed third input makes the call-count comparison use the same primitive
width throughout. An ordinary $2n$-to-$n$ node primitive may instead be used
in both schemes. The number of saved leaf calls stays the same, but the
relative costs of the different primitive widths must then be accounted for.

\subsection{Algorithms}
\paragraph{Commit.}
Compute the leaf digests $v_{0,i}=L_8(X_i)$ for $0\leq i<N$. For
$\ell=0,\ldots,d-1$, compute
\begin{equation}
 v_{\ell+1,j}=C(v_{\ell,2j},v_{\ell,2j+1}),
 \qquad 0\leq j<N/2^{\ell+1}.
 \label{eq:commit}
\end{equation}
The commitment is $R=v_{d,0}$. Store the records and tree digests for later
openings. When $N=1$, there are no internal nodes and the root is simply
the sole leaf digest.

\paragraph{Open.}
To open position $i$, return $X_i$ and the $d$ sibling digests
$\pi=(s_0,\ldots,s_{d-1})$ on its path. Specifically, at level $\ell$ the
current node has index $\lfloor i/2^\ell\rfloor$, and $s_\ell$ is the
other child of its parent. With a stored tree, generating the proof reads
these digests and needs no new compression calls.

\paragraph{Verify.}
Reject unless $0\leq i<N$, the record has eight $n$-bit blocks, and the
proof contains exactly $d$ $n$-bit digests. Set $y_0=L_8(X_i)$. For each
$\ell=0,\ldots,d-1$, let $b_\ell=\lfloor i/2^\ell\rfloor\bmod2$ and set
\begin{equation}
 y_{\ell+1}=
 \begin{cases}
 C(y_\ell,s_\ell),&b_\ell=0,\\
 C(s_\ell,y_\ell),&b_\ell=1.
 \end{cases}
 \label{eq:verify}
\end{equation}
Accept exactly when $y_d=R$. Correctness follows by induction on the level:
the verifier starts with the honest leaf digest and reconstructs the honest
parent at each step.

\paragraph{Communication.}
The authentication proof contains $dn$ bits, and the record contains $8n$
bits. Thus record plus proof costs $(d+8)n$ bits, excluding index and framing.
An explicit fixed-width index adds $d$ bits in either scheme. No extra value
from the computation of $\Tate$ is transmitted. The proof format is the
ordinary binary path; only the rule for hashing the disclosed record changes.

\section{Security}
\label{sec:security}
We use position binding: an adversary should not be able to produce a root
$R$, an index $i$, two distinct records $Z\ne Z'$, and proofs $\pi,\pi'$
that are both accepted at position $i$ under $R$. The root is chosen by the
adversary and need not have a known full opening. This is the relevant
consistency guarantee when an untrusted prover creates the commitment.

\begin{proposition}[Binding]
\label{prop:binding}
For any deterministic leaf function $L:(\{0,1\}^n)^8\to\{0,1\}^n$ and
ordered node function $C$, two accepted openings to different records at the
same position yield either a collision in $L$ or a collision in $C$.
\end{proposition}
\begin{proof}
Evaluate both verification computations, obtaining
$y_0,\ldots,y_d$ and $y'_0,\ldots,y'_d$. If $y_0=y'_0$, then
$L(Z)=L(Z')$ with $Z\ne Z'$, which is a leaf collision.

Otherwise $y_0\ne y'_0$, while acceptance implies $y_d=y'_d=R$.
There is therefore a first level at which the paths merge: for some $\ell$,
$y_\ell\ne y'_\ell$ but $y_{\ell+1}=y'_{\ell+1}$. Because the two
openings use the same index, their current digests occupy the same input
coordinate of $C$. The two ordered input pairs are different and have equal
outputs, giving a collision in $C$. If $d=0$, acceptance directly gives the
leaf-collision case.
\end{proof}

The reduction evaluates only the two supplied paths. It neither constructs
the entire tree nor guesses the level at which a collision occurs. If an
adversary makes $q$ primitive queries before outputting its attack, evaluating
the two openings in our construction adds at most $2(d+3)$ queries. With
$Q=q+2(d+3)$, the resulting bound is
\begin{equation}
 \mtadv_{\mathrm{bind}}(q)
 \leq\epsilon_8(Q)+\epsilon_C(Q),
 \label{eq:binding}
\end{equation}
where $\epsilon_8$ and $\epsilon_C$ bound the probability that a transcript
of at most $Q$ primitive queries contains a fully evaluated collision in
the complete leaf function or the node hash, respectively, with the other
independent oracle interfaces available or simulated. There is no
multiplicative loss in $N$. For an ideal $h_{\mathrm N}$, its restriction
in~\eqref{eq:node} has the usual birthday bound
$\epsilon_C(Q)\leq Q(Q-1)/2^{n+1}$, capped at one.

For $\Tate$, the larger-input extension in~\cite{DKMN21} gives
birthday-scale collision security in the independent ideal-function model.
Section~\ref{sec:collision-bound} gives an explicit refinement of its
collision bound and hence a concrete choice of $\epsilon_8$.
For a concrete instantiation, collision resistance of the complete leaf
function is the property required by Proposition~\ref{prop:binding}.
Collision resistance of the individual compression functions alone is not
a generic proof for this XOR-based composition.

The same argument covers consistency with an honestly committed vector:
pair an alleged false opening with the honest opening at that position.
Likewise, two different vectors with the same root yield a collision by
choosing a differing position and comparing their honest paths. Constructing
any missing honest tree must be charged to the reduction's budget in these
variants.

Crucially, the verifier evaluates all of $\Tate$ on the complete record.
It does not accept an intermediate digest in place of any undisclosed block.
This avoids the additional security questions arising from specialized
partial openings. Position binding is still only a consistency property:
it does not establish hiding, knowledge of a full vector, or the additional
properties that a particular succinct-argument compiler might require.

\subsection{A sharper collision bound for T5 and T8}
\label{sec:collision-bound}
The T5 analysis in~\cite[Appendix~A]{DKMN21} gives the bound
$(n^2+10)q^2/2^n$ by imposing a cutoff of $n$ on certain matching-pair
sets and a cutoff of $nq$ on the number of completed evaluations.
We refine this argument by retaining the factorial in the matching-pair
tail bound, choosing the cutoff freely, and bounding the number of
evaluations in expectation. The proof below treats arbitrary adaptive
query schedules and applies to both T5 and its eight-input extension.

Let $\mathcal{G}=\{0,1\}^n$, with addition written as $\oplus$.
For $\ell\in\{2,3\}$, let $h_1,h_2,h_3:\mathcal{G}^{\ell}\to
\mathcal{G}$ be independent uniformly random functions, and define
\begin{equation}
 H_{\ell}(u,v,s,w)
 =h_3\bigl(h_1(u)\oplus s,h_2(v)\oplus s,w\bigr)\oplus s,
 \label{eq:generalized-t5}
\end{equation}
where $u,v\in\mathcal{G}^{\ell}$, $s\in\mathcal{G}$, and
$w\in\mathcal{G}^{\ell-2}$. For $\ell=2$, the tuple $w$ is empty
and $H_2=\mathsf{T}_5$; for $\ell=3$, $H_3=\Tate$ as defined
in~\eqref{eq:t8}. A transcript \emph{contains an evaluation} of
$H_{\ell}$ if it contains all three primitive queries needed to compute
that evaluation, irrespective of their order. It contains a collision if
two such evaluations have distinct complete inputs and the same output.
We assume throughout this subsection that both claimed collision inputs
have been fully evaluated in the transcript.

\begin{theorem}[Refined collision bound]
\label{thm:refined-collision}
Let $n,Q\geq1$ be integers with $Q\leq2^{n/2}$, and let
$\ell\in\{2,3\}$. Consider any
possibly randomized, computationally unbounded adversary making at most
$Q$ adaptive queries in total to the independent ideal functions
$h_1,h_2,h_3$ above. Set $\rho=Q^2/2^n$ and, for $k\geq1$, define
\begin{equation}
 c_k
 =1+\frac{k^2}{2}+\frac{k\rho}{48}
   +\frac{3\cdot2^n\rho^k}{2^{k+1}(k+1)!}.
 \label{eq:collision-coefficient}
\end{equation}
Starting at $k=1$, stop at the first $k$ for which $c_{k+1}\geq c_k$.
This determines the minimizing cutoff and its coefficient:
\[
 K_\star=\min\{k\in\{1,\ldots,n\}:c_{k+1}\geq c_k\},
 \qquad \gamma_n(Q)=c_{K_\star}.
\]
The cutoff always exists and minimizes $c_k$ over all integers $k\geq1$.
The collision probability satisfies
\begin{equation}
 \Pr[\text{the transcript contains a collision in $H_\ell$}]
 \leq\min\left\{1,\;\gamma_n(Q)\frac{Q^2}{2^n}\right\}.
 \label{eq:refined-collision-bound}
\end{equation}
\end{theorem}

For $n=256$, this gives the following concrete bounds. Each row applies
to every $Q$ up to the indicated query budget; $K_\star$ is the
minimizing cutoff at that budget, and the coefficient is rounded upward.
\begin{center}
\begin{tabular}{@{}rrl@{}}
\toprule
Query budget & $K_\star$ at the budget & Collision probability at most\\
\midrule
$Q\leq2^{64}$  & $2$  & $3.07\,Q^2/2^{256}$\\
$Q\leq2^{80}$  & $3$  & $5.51\,Q^2/2^{256}$\\
$Q\leq2^{96}$  & $4$  & $9.01\,Q^2/2^{256}$\\
$Q\leq2^{112}$ & $8$  & $33.01\,Q^2/2^{256}$\\
$Q\leq2^{120}$ & $13$ & $86.10\,Q^2/2^{256}$\\
\bottomrule
\end{tabular}
\end{center}
For comparison, the original coefficient $n^2+10$ is $65{,}546$.
For $Q>2^{n/2}$ we use the trivial probability bound of one.

\begin{proof}
We first prove the following intermediate bound for every integer
$K\geq1$, and then show that the cutoff in the theorem minimizes it:
\begin{equation}
 B_n(Q,K):=c_K\rho
 =\left(1+\frac{K^2}{2}\right)\rho
  +\frac{K\rho^2}{48}
  +\frac{3\cdot2^n}{(K+1)!}
       \left(\frac{\rho}{2}\right)^{K+1}.
 \label{eq:cutoff-collision-bound}
\end{equation}
Write $S=2^n$ for the size of the output space; the symbol $N$ continues
to denote the number of records in the Merkle tree. Repeated queries can
be answered from a cache, so we count only fresh queries. Pad an execution
with idle rounds if it uses fewer than $Q$ such queries. Sample the
adversary's random tape at the start, and let $\mathcal{F}_i$ be this
tape together with the full history after round $i$. The query in round
$i$ is determined by
$\mathcal{F}_{i-1}$, and each fresh oracle answer is uniform in
$\mathcal{G}$ conditional on this history. All probability and
expectation estimates below use this conditional distribution.

\paragraph{Transcripts and complete evaluations.}
After round $i$, let $\mathcal{A}_i$ contain the query--answer pairs
$(u,a)$ for $h_1$, let $\mathcal{B}_i$ contain $(v,b)$ for $h_2$, and
let $\mathcal{C}_i$ contain $((c,d,w),e)$ for $h_3$. The full input,
including $w$, is part of the identity of an $h_3$ query. For each entry
of $\mathcal{C}_i$, write
\[
 r=e\oplus c,\qquad t=e\oplus d.
\]
A triple consisting of one entry of each table defines a complete
evaluation exactly when
\begin{equation}
 a\oplus b=c\oplus d.
 \label{eq:t5-consistent-triple}
\end{equation}
Indeed, the corresponding input and output are then uniquely determined:
\[
 (u,v,s,w),\qquad s=a\oplus c=b\oplus d,
 \qquad f=a\oplus r=b\oplus t.
\]
Conversely, every evaluation contained in the transcript determines one
such triple. Distinct consistent triples therefore represent distinct
complete inputs. Let $\mathcal{E}_i$ be the set of these evaluations.

\paragraph{Simple collisions.}
Let $\mathsf{Bad}_0$ be the event that two distinct queries have equal
$a$ values in $\mathcal{A}_Q$, equal $b$ values in $\mathcal{B}_Q$,
equal $r$ values in $\mathcal{C}_Q$, or equal $t$ values in
$\mathcal{C}_Q$. Conditional on any history, a fresh $h_1$ or $h_2$
answer collides with a previous answer of that oracle with probability
at most $(i-1)/S$ in round $i$. For $h_3$, each of $e\oplus c$ and
$e\oplus d$ is uniform conditional on the chosen input and history.
A union bound gives probability at most $2(i-1)/S$ for either type
of shifted collision. Independence between these two shifted values
is not needed. Consequently,
\begin{equation}
 \Pr[\mathsf{Bad}_0]
 \leq\sum_{i=1}^{Q}\frac{2(i-1)}{S}
 =\frac{Q(Q-1)}{S}\leq\rho.
 \label{eq:t5-simple-collisions}
\end{equation}

\paragraph{Matching-pair sets and their tails.}
For $z,f\in\mathcal{G}$ define
\begin{align*}
 \mathcal{P}_{12,i}(z)
 &=\{((u,a),(v,b))\in\mathcal{A}_i\times\mathcal{B}_i:
          a\oplus b=z\},\\
 \mathcal{P}_{13,i}(f)
 &=\{((u,a),((c,d,w),e))\in\mathcal{A}_i\times\mathcal{C}_i:
          a\oplus e\oplus c=f\},\\
 \mathcal{P}_{23,i}(f)
 &=\{((v,b),((c,d,w),e))\in\mathcal{B}_i\times\mathcal{C}_i:
          b\oplus e\oplus d=f\}.
\end{align*}
Let $\mathsf{Bad}_K$ be the event that at least one of these sets has
more than $K$ entries after round $Q$. We bound its probability when
$\mathsf{Bad}_0$ does not occur.

Fix one of the three families and one bucket value. In the absence of
simple collisions, a query adds at most one pair to this bucket: the
values in the opposite table are distinct. This also holds for the
shifted tables used in $\mathcal{P}_{13,i}$ and
$\mathcal{P}_{23,i}$. Let $I_i$ indicate that the fixed bucket gains
at least one pair in round $i$, whether or not a simple collision occurs
elsewhere. There are at most $i-1$ possible matching values for the
fresh answer, so for every history
\[
 \Pr[I_i=1\mid\mathcal{F}_{i-1}]\leq\frac{i-1}{S}.
\]
For fixed indices $i_1<\cdots<i_k$, successive conditioning therefore
gives
\[
 \Pr[I_{i_1}=\cdots=I_{i_k}=1]
 \leq\prod_{j=1}^{k}\frac{i_j-1}{S}.
\]
If the bucket has at least $k$ entries and $\mathsf{Bad}_0$ has not
occurred, at least $k$ distinct rounds have $I_i=1$. Taking a union
bound over these rounds and expanding a power yields
\begin{align*}
 \Pr[\text{fixed bucket has at least $k$ entries}
          \;\wedge\;\neg\mathsf{Bad}_0]
 &\leq\sum_{1\leq i_1<\cdots<i_k\leq Q}
             \prod_{j=1}^{k}\frac{i_j-1}{S}\\
 &\leq\frac{1}{k!}
          \left(\sum_{i=1}^{Q}\frac{i-1}{S}\right)^k\\
 &\leq\frac{1}{k!}\left(\frac{\rho}{2}\right)^k.
\end{align*}
The factor $1/k!$ follows because every product with distinct indices
occurs $k!$ times in the power, whose remaining terms are nonnegative.
The argument also covers $k>Q$, when the first sum is empty. There are
$S$ bucket values in each of three families. Setting $k=K+1$ gives
\begin{equation}
 \Pr[\mathsf{Bad}_K\wedge\neg\mathsf{Bad}_0]
 \leq\frac{3S}{(K+1)!}
            \left(\frac{\rho}{2}\right)^{K+1}.
 \label{eq:t5-bucket-tail}
\end{equation}

\paragraph{Expected number of evaluations before a bad event.}
Let $G_i$ be the event that neither a simple collision nor a bucket of
size greater than $K$ has appeared in the first $i$ rounds. Define
\[
 V_i=\mathbf{1}_{G_i}|\mathcal{E}_i|,
 \qquad V_0=0.
\]
Thus the count is set to zero once a bad event occurs. We bound
$\mathbb{E}[V_i]$ without conditioning on the future event $G_Q$.

Every evaluation is first completed by exactly one query. Let $L_i$
count evaluations first completed by a query to $h_1$ or $h_2$ in the
first $i$ rounds of the original execution, including executions on
which a bad event occurs. For a fresh $h_1$ query in round $j$, each
pair in $\mathcal{B}_{j-1}\times\mathcal{C}_{j-1}$ completes an
evaluation precisely when its required value $b\oplus c\oplus d$
equals the fresh answer $a$. Linearity of expectation gives conditional
expected increment
\[
 \frac{|\mathcal{B}_{j-1}|\,|\mathcal{C}_{j-1}|}{S}
 \leq\frac{(j-1)^2}{4S}.
\]
Here the two table sizes sum to at most $j-1$. The same bound applies
to a fresh $h_2$ query, with $\mathcal{A}_{j-1}$ in place of
$\mathcal{B}_{j-1}$. An $h_3$ query or an idle round does not increase
$L_i$. In particular, these estimates hold for adaptive, random
allocations of queries among the oracles, and
\begin{equation}
 \mathbb{E}[L_i]\leq\frac{1}{4S}\sum_{j=1}^{i}(j-1)^2.
 \label{eq:t5-late-inner-evaluations}
\end{equation}

On $G_i$, each $h_3$ query with input $(c,d,w)$ completes exactly one
evaluation per entry of $\mathcal{P}_{12,j-1}(c\oplus d)$ at its
query time $j$, hence at most $K$ evaluations. There are at most $i$
such queries. Thus $V_i\leq Ki+L_i$ on every execution: this is the
preceding counting argument on $G_i$, and is immediate from $V_i=0$
otherwise. We obtain
\begin{equation}
 \mathbb{E}[V_i]
 \leq Ki+\frac{1}{4S}\sum_{j=1}^{i}(j-1)^2.
 \label{eq:t5-expected-evaluations}
\end{equation}

\paragraph{Collisions created by a fresh query.}
Fix a history satisfying $G_{i-1}$. First, two distinct evaluations
completed by the \emph{same} query cannot collide. For an $h_1$ query,
their common new value $a$ and equal outputs would force $r=r'$.
The absence of shifted collisions identifies the two $h_3$ queries.
Consistency then gives $b=b'$, and the absence of collisions in $h_2$
identifies the two $h_2$ queries as well. The triples would be identical.
The $h_2$ case is symmetric, using $t=t'$. For an $h_3$ query, the
common $r$ and equal outputs force $a=a'$, and consistency then forces
$b=b'$. Again, the absence of simple collisions in the previous tables
identifies both triples. These arguments use only the good prefix
$G_{i-1}$.

It remains to bound collisions between a new evaluation and an old
evaluation in $\mathcal{E}_{i-1}$. Fix one old evaluation, with output
$f$. There are three cases.
\begin{enumerate}
 \item A fresh query to $h_1$ returns $a$. A new evaluation with output
 $f$ must use a pair in $\mathcal{P}_{23,i-1}(f)$ and must satisfy
 $a=b\oplus c\oplus d$. There are at most $K$ such pairs, each
 specifying one value of the uniform answer $a$. The probability is
 at most $K/S$.
 \item A fresh query to $h_2$ returns $b$. The pair must lie in
 $\mathcal{P}_{13,i-1}(f)$, and its required answer is
 $b=a\oplus c\oplus d$. The same bound $K/S$ follows.
 \item A fresh query to $h_3$ has its input $(c,d,w)$ fixed before
 the answer $e$ is sampled. A new evaluation uses a pair in
 $\mathcal{P}_{12,i-1}(c\oplus d)$ and has output $f$ precisely
 when $e=f\oplus a\oplus c$. Again, there are at most $K$
 candidate answers, so the probability is at most $K/S$.
\end{enumerate}
Only one of these cases applies in each round; their bounds are not
added. A union bound over the old evaluations shows that, conditional
on the fixed good history, the probability of a newly created collision
is at most $K|\mathcal{E}_{i-1}|/S$. Equivalently, if $X_i$ is the
event that a collision involving a newly completed evaluation appears
in round $i$, then
\begin{equation}
 \Pr[X_i\wedge G_{i-1}]
 \leq\frac{K}{S}\mathbb{E}[V_{i-1}].
 \label{eq:t5-per-query-collision}
\end{equation}

\paragraph{Summing the bounds.}
If a collision is present at the end and $G_Q$ holds, it was first
created in some round whose prefix was good. Therefore,
\begin{align*}
 \Pr[\text{collision}\wedge G_Q]
 &\leq\sum_{i=1}^{Q}\frac{K}{S}\mathbb{E}[V_{i-1}]\\
 &\leq\frac{K^2}{S}\sum_{i=1}^{Q}(i-1)
   +\frac{K}{4S^2}\sum_{i=1}^{Q}\sum_{j=1}^{i-1}(j-1)^2\\
 &=\frac{K^2Q(Q-1)}{2S}
   +\frac{KQ(Q-1)^2(Q-2)}{48S^2}\\
 &\leq\frac{K^2\rho}{2}+\frac{K\rho^2}{48}.
\end{align*}
For $Q=1$ both sums and both terms in the penultimate line are zero.
For $Q\geq2$, the double-sum identity used here is
\[
 \sum_{i=1}^{Q}\sum_{j=1}^{i-1}(j-1)^2
 =\sum_{r=0}^{Q-2}(Q-1-r)r^2
 =\frac{Q(Q-1)^2(Q-2)}{12}.
\]
Finally, decompose the collision event according to
$\mathsf{Bad}_0$, $\mathsf{Bad}_K\wedge\neg\mathsf{Bad}_0$, and
$G_Q$. Adding~\eqref{eq:t5-simple-collisions},
\eqref{eq:t5-bucket-tail}, and the preceding collision estimate gives
$B_n(Q,K)=c_K\rho$ for every $K\geq1$.

\paragraph{The minimizing cutoff.}
To justify the explicit stopping rule, write
\[
 a_k=\frac{3\cdot2^n\rho^k}{2^{k+1}(k+1)!},
 \qquad c_k=1+\frac{k^2}{2}+\frac{k\rho}{48}+a_k.
\]
The tail coefficients obey the recurrence
$a_{k+1}=\rho a_k/(2(k+2))$. Consequently,
\begin{align*}
 c_{k+1}-c_k
 &=k+\frac12+\frac{\rho}{48}
       -a_k\left(1-\frac{\rho}{2(k+2)}\right),\\
 c_{k+2}-2c_{k+1}+c_k
 &=1+a_k\left(1-\frac{\rho}{k+2}
                 +\frac{\rho^2}{4(k+2)(k+3)}\right)>0.
\end{align*}
The last inequality uses $0<\rho\leq1$ and $k\geq1$. Thus the
successive differences are strictly increasing: once the coefficients
stop decreasing, they never decrease again. Moreover,
$a_n\leq3/(2(n+1)!)$, so
\[
 c_{n+1}-c_n
 \geq n+\frac12-\frac{3}{2(n+1)!}>0.
\]
Hence the stopping rule terminates by $k=n$ and returns the smallest
integer attaining the global minimum. Substituting $K=K_\star$ into
the intermediate bound and capping the result at one proves the theorem.
Finally, each $c_k$ is nondecreasing in $\rho$. A coefficient evaluated
at a query-budget endpoint therefore also bounds all smaller query
counts, which justifies the ranges in the table.
\end{proof}

\paragraph{Uniform asymptotic bound.}
The optimized coefficient also admits a simple asymptotic estimate.
For any $K$ satisfying
\[
 2^{K+1}(K+1)!\geq3\cdot2^n,
\]
the last term of~\eqref{eq:cutoff-collision-bound} is at most
$\rho^{K+1}\leq\rho$. Since $\gamma_n(Q)\leq c_K$, we obtain
\begin{equation}
 \gamma_n(Q)\leq2+\frac{K^2}{2}+\frac{K}{48}.
 \label{eq:t5-uniform-cutoff}
\end{equation}
Stirling's formula permits $K=O(n/\log n)$ as $n\to\infty$, giving
an $O((n/\log n)^2Q^2/2^n)$ collision bound in this range.
This improvement is an upper-bound refinement; we do not claim that the
optimized bound is tight or that an absolute constant independent of $n$
suffices.

\paragraph{Consequence for the Merkle commitment.}
With $Q=q+2(d+3)\leq2^{n/2}$ as in~\eqref{eq:binding}, the theorem
and the node-hash birthday bound give
\begin{equation}
 \mtadv_{\mathrm{bind}}(q)
 \leq\min\left\{1,\;
 \left(\gamma_n(Q)+\frac12\right)\frac{Q^2}{2^n}\right\}.
 \label{eq:refined-binding}
\end{equation}
All evaluations of the two verification paths are already charged to
$Q$. Other independent oracle interfaces can be simulated with
independent randomness in the
leaf-collision analysis. The proof and the bound therefore remain valid
when the adversary also uses the domain-separated node and baseline
interfaces of our construction.

\section{Comparison with a standard Merkle tree}
\label{sec:cost}
\subsection{Baseline and exact counts}
Use the same records, digest length, and binary upper tree. Hash a baseline
leaf by fixed-length Merkle--Damg\aa rd iteration:
\begin{equation}
\begin{aligned}
 z_0&=\mathsf{IV},\\
 z_j&=h_{\mathrm B}(z_{j-1},x_{2j-1},x_{2j})
      \quad(j=1,2,3,4),\qquad L_{\mathrm{MD}}(X)=z_4.
\end{aligned}
\label{eq:baseline}
\end{equation}
The public initialization value and input length are fixed, so no extra
padding call is charged. Each leaf uses four calls to a $3n$-to-$n$
compression function. This is the particular standard leaf-hashing baseline
for the quantitative comparison, rather than a claim about every hash API.

A binary tree with $N$ leaves has $N-1$ internal nodes. Commitment therefore
uses $4N+(N-1)$ baseline calls and $3N+(N-1)$ calls with $\Tate$.
Verification evaluates one leaf hash and $d$ node hashes. These observations
give the following exact comparison.

\begin{center}
\begin{tabular}{@{}lcc@{}}
\toprule
Quantity & Standard leaves & $\Tate$ leaves\\
\midrule
Calls to commit & $5N-1$ & $4N-1$\\
Calls to verify one record & $d+4$ & $d+3$\\
Calls to replace a complete record & $d+4$ & $d+3$\\
Authentication proof & $dn$ bits & $dn$ bits\\
Opened record & $8n$ bits & $8n$ bits\\
Root & $n$ bits & $n$ bits\\
Stored tree digests & $2N-1$ & $2N-1$\\
\bottomrule
\end{tabular}
\end{center}

The update row assumes that the upper tree is stored and that the new leaf
hash is recomputed from scratch. The committer saves $N$ calls overall, and
the verifier saves one call per opened record. The relative reductions are
\begin{equation}
 \rho_{\mathrm{commit}}=\frac{N}{5N-1},\qquad
 \rho_{\mathrm{verify}}=\frac{1}{d+4}.
 \label{eq:ratios}
\end{equation}
Thus leaf hashing uses $25\%$ fewer calls, while total commitment work
approaches a $20\%$ reduction. The verification saving is additive and its
percentage decreases as the tree grows. Under an equal-cost,
compression-dominated model, the commitment ratio corresponds to a limiting
throughput speedup of $5/4$.

\subsection{Concrete examples}
For $n=256$, each record is 256 bytes and each sibling digest is 32 bytes.
The following counts exclude any common context binding, index, or framing.

\begin{center}
\small
\begin{tabular}{@{}rrrr@{}}
\toprule
$N$ & Commit: standard / $\Tate$ & Verify: standard / $\Tate$ & Path bytes\\
\midrule
$2^{10}$ & $5{,}119\;/\;4{,}095$ & $14\;/\;13$ & $320$\\
$2^{16}$ & $327{,}679\;/\;262{,}143$ & $20\;/\;19$ & $512$\\
$2^{20}$ & $5{,}242{,}879\;/\;4{,}194{,}303$ & $24\;/\;23$ & $640$\\
\bottomrule
\end{tabular}
\end{center}

At $N=2^{20}$, the committer saves $1{,}048{,}576$ calls and verification
uses 23 rather than 24 calls, a $4.17\%$ reduction. Record plus path is
896 bytes in both cases. These numbers illustrate the computational gain
without conflating authentication digests with the record being opened.

\subsection{Compression counts versus running time}
Suppose a leaf-compression call costs $\tau$, an internal-node call costs
$\nu$, and an $n$-bit XOR costs $\xi$. Assuming the specialized and baseline
compression roles have the same cost, the work estimates are
\begin{align}
 T_{\mathrm{standard}}&=4N\tau+(N-1)\nu,\\
 T_{8}&=3N\tau+3N\xi+(N-1)\nu.
\end{align}
The saved work is $N(\tau-3\xi)$, before accounting for any difference
in domain-separation overhead. If the primitive roles have unequal costs,
replace $3\tau$ by their sum. Memory traffic, batching, and instruction-level
parallelism can also affect elapsed time. Accordingly, a call-count saving
is a useful model rather than an implementation benchmark.

The first two calls of $\Tate$ can run concurrently. With sufficient
processors, the compression critical path from the raw records to the root
has length $d+2$, compared with $d+4$ for the sequential MD baseline. A
conventional tree-shaped leaf hash can itself exploit parallelism, so this
depth comparison is specific to the baseline in~\eqref{eq:baseline}.

\section{Batched openings and updates}
\label{sec:batch}
\paragraph{Multiproofs.}
Let $S$ contain $u$ distinct record positions. A standard Merkle multiproof
supplies the boundary digests needed to reconstruct the union of their paths.
Let $B(S)$ be the number of transmitted boundary digests and $K(S)$ the
number of internal hashes evaluated during reconstruction. The tree shape
and queried positions are unchanged by our leaf substitution, so the
verification costs are
\begin{equation}
 4u+K(S)\quad\text{versus}\quad3u+K(S),
 \label{eq:multiproof}
\end{equation}
and the authentication data has size $B(S)n$ in both cases. The saving is
one call per distinct record whose digest is computed, not per repeated
query. Existing path deduplication and leaf caching can be applied equally
to both constructions.

For example, opening two sibling leaves at positions $2j$ and $2j+1$ in a
tree with $d\geq1$ requires
$d-1$ boundary digests: the verifier computes their common parent and follows
its path. It performs $d$ internal hashes and either eight or six leaf calls.
Thus verification drops from $d+8$ to $d+6$, with unchanged communication.
At the other extreme, opening every record uses no boundary digests and
reduces to the full commitment comparison.

\paragraph{Updates and storage.}
Replacing a whole record requires its new leaf digest and the $d$ ancestors
on its path. It saves one compression call under the same comparison.
If intermediate leaf values are cached and only one block changes, the
counts depend on its position. In $\Tate$, changing a block among the first
six needs two leaf calls, while changing either of the last two needs only
the final call. A cached MD chain also admits cheaper suffix recomputation.
The uniform one-call saving therefore concerns whole-record replacement.

Both schemes store the same $2N-1$ tree digests and the same records.
Caching $\Tate$ intermediates is optional and adds memory beyond this count.
An old commitment cannot be converted by changing its root alone: new leaf
digests and their ancestors must be computed under the new hash suite.

\paragraph{Other tree shapes.}
For a public non-power-of-two record count, use an agreed binary tree shape
or canonical padding. An unpadded full binary tree still has $N-1$ internal
nodes, although path lengths can differ. For a complete $a$-ary upper tree,
with one primitive call per internal node, the commitment counts become
\[
 4N+\frac{N-1}{a-1}
 \quad\text{and}\quad
 3N+\frac{N-1}{a-1}.
\]
The authentication proof has $(a-1)\log_aN$ digests in either construction.
The leaf substitution preserves proof size within each chosen tree shape;
it does not make binary and higher-arity proofs equal in size. If wider
nodes require more than one call, their common cost must be included.

\paragraph{Longer records.}
The same approach extends to some larger fixed formats while retaining a
complete-record opening. For a record of $7k+1$ blocks, first hash eight
blocks with $\Tate$, then hash the resulting digest together with the next
seven blocks, repeating until all blocks have been consumed. More precisely,
\[
 s_1=\Tate(x_1,\ldots,x_8),\qquad
 s_j=\Tate(s_{j-1},x_{7j-5},\ldots,x_{7j+1})
 \quad(2\leq j\leq k).
\]
The leaf digest is $s_k$. This uses $3k$ compression calls, compared with
$\lceil(7k+1)/2\rceil$ for the fixed-length MD baseline, padding its final
input with a fixed word when necessary. For example, a fifteen-block record
uses six rather than eight calls. The original eight-block construction is
the case $k=1$; the upper Merkle tree remains unchanged.

For a fixed record length, a collision in this iterated leaf hash gives a
$\Tate$ collision: trace the equal final digests backward until the complete
input tuples of two corresponding invocations differ. If their tuples
never differ, the records themselves are equal. The full record is still
disclosed in an opening, so the tree-level binding reduction applies as
before. Different record lengths require an unambiguous encoding and
authenticated context. This extension is useful only when those longer
records are already the application's opening unit; it does not justify
enlarging records without accounting for the resulting communication.

\section{Concrete compression interfaces and schedules}
\label{sec:concrete}
The Rust implementation in \path{impl/} specializes the Plonky3 binary
Merkle construction and provides an optimized standard-hash comparator.
Our workload contains $M=2^{20},\ldots,2^{26}$ 32-bit field encodings,
with $W=2^0,\ldots,2^{14}$ elements per record. Thus a leaf contains
$B=4W$ bytes and the tree has $N=M/W$ leaves. The implementation treats
these as little-endian u32 encodings; any field-canonicality check belongs
to the surrounding protocol. In this section digests have $n=256$ bits,
or 32 bytes. Every opening still reveals the entire record.

\subsection{Count the actual primitive, including padding}
An abstract $3n$-to-$n$ call need not cost one native compression. In our
tagged SHA-256 adapter, a 17-byte tag followed by 96 input bytes occupies
two padded SHA-256 blocks. Consequently, ordinary T8 costs six native
calls for a 256-byte leaf, whereas standard SHA-256 costs five. Tagged
SHA3-256 can hash the same 113 bytes in one permutation, but standard
SHA3-256 hashes a 256-byte leaf in only two permutations, fewer than
T8's three. This explains why merely translating the abstract gadget can
lose even before accounting for XORs, packing, and function dispatch.

For a nonempty $B$-byte leaf, the standard implementations use
\begin{align}
 c_{\mathrm{SHA256}}(B)
   &=\floor{B/64}+1+\mathbf{1}[B\bmod64\geq56],\nonumber\\
 c_{\mathrm{SHA3}}(B)&=\floor{B/136}+1,\nonumber\\
 c_{\mathrm{BLAKE3}}(B)
   &=\ceil{B/64}+\ceil{B/1024}-1.
 \label{eq:native-standard}
\end{align}
The last expression includes BLAKE3's internal chunk-tree compressions.
One unit means one SHA-256 or BLAKE3 compression, or one full 24-round
Keccak-f[1600] permutation. Their execution times are different. Each
backend's unchanged binary parent costs one such unit, giving
\[
 C_{\mathrm{commit}}=N c(B)+N-1,\qquad
 C_{\mathrm{verify}}=c(B)+\log_2N.
\]
These native counts are distinct from the fixed-format MD comparator in
\eqref{eq:baseline}. Padding boundaries can remove a saving at some widths;
there is no strict improvement at every parameter size.

\subsection{Use the complete native input while reserving real role bits}
\paragraph{SHA-256: 95 payload bytes per call.}
A raw SHA-256 compression accepts a 32-byte chaining value (CV) and a
64-byte message block. Reserve one CV byte for a role $\rho$ and use the
remaining 31 CV bytes as payload. For a 95-byte string $u$, define
\[
 h_\rho(u)=\mathsf{Compress}_{\mathrm{SHA256}}
       (\rho\concat u[64{:}95],\;u[0{:}64]).
\]
Here slices have byte offsets and exclude their right endpoint; native SHA
words are decoded and the result encoded in big-endian order. Distinct role
bytes specify disjoint subsets of the raw compressor's input space. XORing
a role constant into an \emph{unrestricted} CV would not achieve this:
each role would still have the same domain. Nor may one discard a byte of
a previous digest just to make room for a tag.

\paragraph{BLAKE3: seven counter bytes are additional payload.}
The native BLAKE3 interface also contains a 64-bit counter~\cite{BLAKE3}.
For a 103-byte payload, use bytes $0{:}64$ for the block, $64{:}96$ for
the CV, and $96{:}103$ for the low 56 counter bits in little-endian order.
The high counter byte is the role. Fix the block length to 64 and flags to
\texttt{KEYED\_HASH | CHUNK\_START | CHUNK\_END} (19), without
\texttt{ROOT}. All 103 payload bytes then enter one compression. These are
research oracles over raw compression, not calls to standard BLAKE3 hashing;
the counter payload is not a sequential chunk index.

\subsection{Widen the gadget and put the dependency at its root}
Let $m\geq64$ be the payload bytes accepted by one role oracle. The
generalized T gadget takes two $m$-byte strings $u,v$, a 32-byte shared
word $s$, and an $(m-64)$-byte extension $e$:
\[
 a=h_0(u),\quad b=h_1(v),\qquad
 T(u,v,s,e)=h_2(a\oplus s,b\oplus s,e)\oplus s.
\]
It consumes $3m-32$ bytes in three calls. In subsequent stages, place the
previous digest in $s$ and absorb $3m-64$ fresh bytes. The first two calls
of every stage then depend only on fresh input, enabling lookahead across
stages. This placement differs from the bottom-input chaining described
in the earlier longer-record example. It changes the leaf function, but
not its input capacity or number of calls.

For ABR, we use only the height-three, seven-call gadget for which the
corrected analysis treats the ordinary narrow case~\cite{DDN22}. Its
widened variant takes four $m$-byte bottom inputs and three pairs
$(s_j,e_j)$ of lengths 32 and $m-64$. With independent roles, compute
\begin{align*}
 a_i&=h_i(u_i) &&(0\leq i<4),\\
 b_j&=h_{4+j}(a_{2j}\oplus s_j,a_{2j+1}\oplus s_j,e_j)
           \oplus a_{2j+1} &&(j=0,1),\\
 A&=h_6(b_0\oplus s_2,b_1\oplus s_2,e_2)\oplus b_1.
\end{align*}
The feed-forward is the right child, whereas T feeds forward its shared
word. The widened ABR gadget consumes $7m-96$ bytes. Chaining through
the \emph{root shared word} $s_2$ leaves $7m-128$ fresh bytes per stage
and makes all six non-root calls independent of the preceding digest.
This observation improves the verification schedule as well as commitment
batching. It does not establish security of the widening or of arbitrary
ABR heights; see Section~\ref{sec:concrete-security}.

\begin{center}
\small
\begin{tabular}{@{}llrrr@{}}
\toprule
Backend & Gadget & First bytes & Later fresh bytes & Calls\\
\midrule
SHA-256 & T253 ($m=95$) & 253 & 221 & 3\\
SHA-256 & ABR-wide ($m=95$) & 569 & 537 & 7\\
BLAKE3 & T277 ($m=103$) & 277 & 245 & 3\\
BLAKE3 & ABR-wide ($m=103$) & 625 & 593 & 7\\
\bottomrule
\end{tabular}
\end{center}
The T253 and T277 names describe first-stage byte capacities; they are
local names for concrete adapters. The older T8 and T253 implementations
retain their original schedules and roots.

\subsection{Choose a short tail instead of paying for a whole gadget}
Let a gadget cost $g$ calls, with first capacity $32+D$ and subsequent
capacity $D$. Thus $(g,D)=(3,3m-64)$ for T and $(7,7m-128)$ for ABR.
A restricted MD call packs the \emph{full} previous 32-byte digest and
$t=m-32$ fresh bytes into its oracle input. Reserve a separate MD role.
With $x_+=\max\{0,x\}$, compare
\begin{align}
 C_0(B)&=1+\ceil{(B-m)_+/t},\nonumber\\
 C_k(B)&=gk+\ceil{(B-32-kD)_+/t}\qquad(k\geq1).
 \label{eq:hybrid-planner}
\end{align}
The first option uses one direct MD head and then tails; the second uses
gadgets followed by tails, or a final zero-padded gadget. The public width
determines the plan. Fixed-width zero padding is injective within that
width; different widths must be distinguished by the authenticated context.
Tie-breaking chooses least padding, then most gadgets.

An exhaustive search over the record length is unnecessary. Among complete
gadgets, \eqref{eq:hybrid-planner} equals
$\ceil{(B-32-k(D-gt))/t}$. Since $D-gt$ is 32 for T and 96 for ABR,
this cost is nonincreasing in $k$. Integer arithmetic finds its minimum
and the first $k$ attaining it, which minimizes padding among those ties.
Compare that plan with the MD-only and final-padded-gadget candidates.
The implementation uses this $O(1)$ planner and tests it against exhaustive
enumeration.

At 64 KiB, SHA-256 ABR-wide pads its 122nd gadget and costs 854 calls.
BLAKE3 ABR-wide uses 110 gadgets and four MD tails, costing 774.
T277 uses 266 gadgets and five tails, costing 803; 267 gadgets and two
tails have the same call count but more padding. New T roles use
\texttt{C0--C2} with MD role \texttt{C3}, and ABR uses
\texttt{D0--D6} with MD role \texttt{D7}, all hexadecimal.
The older SHA-256 T253 mode instead has an unrestricted 64-byte raw-MD
tail. Its domain overlaps the restricted gadget domains, so it is not the
same role-separated hybrid and needs a separate composition argument.

\subsection{For Keccak, change the absorption mode}
SHA3-256 already absorbs 136 bytes per permutation. A chained T gadget
absorbs $m-64/3$ fresh bytes per call, so even filling an oracle to
$m=136$ bytes cannot beat that rate. A useful adapter would need
$m>157\frac13$; the usual tagged SHA3 wrapper cannot provide it.

SHAKE128 has a 168-byte rate~\cite{FIPS202}. Taking its first 32 output
bytes gives the exact standard XOF, with no extra prefix and
$\floor{B/168}+1$ permutations. It targets 128-bit classical preimage
security, unlike SHA3-256's 256-bit preimage target. This is an explicit
choice of a different leaf hash, not an implementation of SHA3-256 with
fewer rounds.

SPONGE-DM~\cite{SPONGEDM} offers another tradeoff. Let $P$ be full-round
Keccak-f[1600], with capacity 272 and rate 1,328 bits (166 bytes).
Starting with a zero 200-byte state, absorb each padded rate block $v$ by
\[
 x=\mathsf{state}\oplus(v\concat0^{272}),\qquad
 \mathsf{state}=P(x)\oplus x.
\]
The feed-forward covers the \emph{entire} state, not only the capacity.
Our encoding prepends the 17 bytes
\texttt{MTLFv001} followed by hexadecimal
\texttt{04 50 00 00 00 00 00 00 00}, and uses $\mathsf{pad10^*1}$.
Its first 32 output bytes need no further permutation. The leaf costs
$\floor{(B+17)/166}+1$ calls. Both Keccak alternatives retain the
ordinary SHA3-256 binary parents.

At $B=65{,}536$, the useful leaf-call counts are as follows. No row implies
that the modes have interchangeable security guarantees.
\begin{center}
\small
\begin{tabular}{@{}llrrr@{}}
\toprule
Backend & Leaf mode & Standard & Variant & Fewer calls\\
\midrule
SHA-256 & ABR-wide & 1,025 & 854 & $16.7\%$\\
BLAKE3 & ABR-wide & 1,087 & 774 & $28.8\%$\\
BLAKE3 & T277 & 1,087 & 803 & $26.1\%$\\
Keccak & SHAKE128 & 482 & 391 & $18.9\%$\\
Keccak & SPONGE-DM272 & 482 & 395 & $18.0\%$\\
\bottomrule
\end{tabular}
\end{center}

\subsection{What is proved, and what remains an assumption}
\label{sec:concrete-security}
Collision resistance of a compression function does not, on its own,
imply collision resistance of T or ABR~\cite{CN24}. Their analyses use
independent ideal role functions. Disjoint role encodings are necessary
for the intended model, but do not prove ideal behavior of concrete raw
SHA-256 or variable-CV, variable-counter BLAKE3 compression. The earlier
T8 theorem is not a security theorem for every adapter in this section.

The corrected ABR result concerns the narrow height-three gadget; it
neither proves arbitrary-height ABR nor validates its specialized partial
openings~\cite{DDN22}. Our widened ABR adaptation has a provisional
classical reduction sketch. For each role, map each distinct wide query
$(x,y,z)$ injectively to a previously unused narrow pair $(u,v)$ with
$u\oplus v=x\oplus y$, without inspecting its output. Each XOR class
has $2^n$ pairs; require fewer than $2^n$ mapped queries per role,
including completion queries, so these assignments cannot exhaust a class.
Writing $(u,v)=(x\oplus r,y\oplus r)$ lets one replace
the shared word of an internal vertex by $s\oplus r$ while preserving
its output and right-child feed-forward. The recorded injective maps then
recover a distinct narrow message pair from a wide collision. Complete
both colliding evaluations first (at most 14 additional queries), and
extract a single-gadget or MD-tail collision before applying this argument
to a chain. This sketch still requires formal review and retains the
polynomial loss of the corrected narrow bound; we do not claim a tight
concrete 128-bit guarantee.

For SPONGE-DM, the ideal-permutation analysis gives generic classical
exponents $\min(n/2,c/2)$ for collisions, $n$ for preimages, and
$\min(n,c-\log_2\alpha)$ for second preimages, subject to the paper's
parameter conditions~\cite{SPONGEDM}. Our equal input/output capacities
$c=c'=272$ satisfy those conditions. Up to 64 KiB, $\alpha\leq395$
absorbed blocks, so these exponents are 128, 256, and 256. This is not
a proof of concrete Keccak security or a guarantee for arbitrary record
lengths. For T5, even an $n$-bit output admits a generic preimage attack
around $2^{2n/3}$ queries~\cite{DKMN21}; collision and preimage claims
must not be conflated.

The same complete-record binding reduction applies once the actual leaf
and parent functions have the needed collision resistance. It does not
establish extractability, quantum security, or substitutability for a
random oracle in a surrounding proof system. Candidate bit-packed role
encodings and the proposed Keccak T520 adapter remain unimplemented; in
particular T520's joint-role composition is unfinished. The repository's
\path{impl/NEXT_CONSTRUCTIONS.md} records these candidates and assumptions.

\section{Engineering and measured performance}
\label{sec:measurements}
\subsection{Reduce packing overhead and expose independent work}
\paragraph{Fuse the whole leaf computation.}
Saving a compression can be outweighed by serializing each intermediate
digest into bytes and packing it back into primitive words. The native
SHA-256 ABR kernel reads 95-byte payloads directly, retains digest and
chaining values in vector registers, and interleaves independent SHA2
instructions. Its two-record commitment path also pairs the otherwise
dependent root and tail calls across records. The BLAKE3 NEON kernel
processes four records, including variable CVs and counters, and keeps
intermediates in vector form across gadgets and MD tails.

For single-opening verification there are no other records to fill lanes.
Instead, bounded lookahead prepares independent branches of four BLAKE3
stages, then evaluates their dependent roots in order. SHA verification
pairs independent calls within a gadget. Scalar and two-record Keccak
paths retain their permutation state across all absorbed blocks, dispatch
once per leaf or batch, and pack full blocks word-wise even at the
unaligned 166-byte rate. These paths use bounded scratch and no per-leaf
heap allocation; runtime feature checks select native or portable kernels.

\paragraph{Distinguish useful calls from occupied SIMD lanes.}
SIMD does not change the mathematical call count, but a partially filled
window can repeat work in unused lanes. For a 64 KiB BLAKE3 leaf, ABR-wide
verification executes 786 compression lanes for 774 useful calls, and
T277 executes 807 for 803. Four-record commitment batches occupy every
lane. Count plots report useful calls, while timing includes this overhead.

\paragraph{Optimize the baseline and the tree schedule as well.}
The same independent-message batching helps standard SHA-256 without
changing any digest or call count: process two leaves with separate IVs,
standard padding, and length encodings through the SHA2 kernel. Standard
BLAKE3 receives complete leaf slices to expose its own chunk SIMD; small
equal-width leaves use its native multi-message interface. Tree digests
are stored contiguously, leaf jobs are coarse enough to preserve SIMD
batches, and small upper levels run serially. Recommit reuses storage,
but fresh-commit measurements include allocation and destruction.
Larger jobs and extra parent batching did not consistently improve the
measured cases. Worker count is a machine-dependent tuning parameter,
not a reason to assume linear scaling.

\subsection{Measurements and the importance of the comparator}
The 28 September 2026 sweep used an Apple M4 Max, four Rayon workers,
Rust 1.96.0, and seven calibrated batch samples per case, targeting
150 ms of sampling. It covered all 15 widths at $M=2^{20}$ and widths
$W\in\{64,256,1024,16384\}$ at $M=2^{24},2^{26}$, for 460 timing rows.
The table below uses $M=2^{26}$ and $W=2^{14}$ (256 MiB total, 64 KiB
per leaf). Each speedup is the same executable's standard median divided
by the variant median for the same backend and geometry.

\begin{center}
\small
\setlength{\tabcolsep}{5pt}
\begin{tabular}{@{}llrrrr@{}}
\toprule
Backend & Mode & Commit ms & Speedup & Verify $\mu$s & Speedup\\
\midrule
SHA-256 & Standard & 25.959 & $1.00\times$ & 22.541 & $1.00\times$\\
 & ABR-wide & 14.300 & $1.82\times$ & 14.300 & $1.58\times$\\
BLAKE3 & Standard & 28.842 & $1.00\times$ & 27.956 & $1.00\times$\\
 & ABR-wide & 20.611 & $1.40\times$ & 25.510 & $1.10\times$\\
 & T277 & 21.433 & $1.35\times$ & 30.073 & $0.93\times$\\
Keccak & SHA3-256 & 34.635 & $1.00\times$ & 65.955 & $1.00\times$\\
 & SHAKE128 & 27.798 & $1.25\times$ & 55.520 & $1.19\times$\\
 & SPONGE-DM272 & 29.274 & $1.18\times$ & 58.245 & $1.13\times$\\
\bottomrule
\end{tabular}
\end{center}
These are historical measurements against that version of the optimized
standard kernel. The subsequent standard SHA-256 pair batching was measured
separately on 29 September on an Apple M3. Across its sampled four-worker
cases, standard commitment improved by $1.12$--$1.46\times$, preserving
roots and leaving verification unchanged. It would be incorrect to divide
the M4 ratios by those M3 improvements to obtain current speedups.

The M3 control sweep does provide a direct comparison within the newer
executable. At $M=2^{20}$ and 64 KiB leaves, standard SHA-256 commitment
took 0.3372 ms versus 0.2926 ms for ABR-wide ($1.15\times$), and
verification took 26.205 versus 15.571 $\mu$s ($1.68\times$).
These are medians of two process medians, each using seven calibrated
samples targeting 180 ms. They illustrate why an improved standard
comparator can substantially reduce a construction's apparent commitment
advantage without reducing its verification advantage: within this same
M3 control sweep, the before executable gave ABR-wide a $1.67\times$
commitment advantage, while its verification advantage was about
$1.68\times$ in both executables.

The improvements are not uniform. In the M4 sweep, BLAKE3 ABR-wide
commitment at $W=16$ was only $0.80\times$ standard speed; at $W=64$
it was $0.85\times$. Short Keccak leaves often have identical permutation
counts, leaving feed-forward and packing overhead exposed. BLAKE3 T277's
large-leaf verification is slower despite fewer useful calls. A local
BLAKE3 speedup peak at 2 KiB records also reflects a standard-kernel
dispatch boundary rather than a uniformly better asymptotic rate.

\paragraph{Reproducibility and validation.}
Backends and before/after processes run sequentially, with builds, input
generation and proof preparation outside timing. Fresh commitment includes
allocation and drop; verification includes a complete leaf and its path.
Case order is fixed, and percentile spreads are not confidence intervals.
The repository retains raw timings, commands, CPU and compiler details,
and source/executable fingerprints in \path{impl/results/}.
\path{impl/LOWCALL_RESULTS.md} describes the M4 sweep;
\path{impl/CODE_SPEED_RESULTS.md} describes the M3 baseline optimization.
Neither sweep establishes an end-to-end SNARK speedup.

Independent scalar references, padding and schedule boundaries, odd SIMD
batches, count-model agreement, and complete-opening tampering tests check
the implementations. The recorded validation passes 76 tests and a doctest
with and without parallelism, plus forced software Keccak checks. These
are correctness checks on the tested architecture, not cryptographic proofs.

\section{Applications to SNARKs and authenticated data}
\label{sec:applications}
\subsection{Merkle-based succinct arguments}
A natural application is an argument whose prover commits to oracle tables
and opens the records queried by its verifier. IOP-based compilation gives
a general setting for such commitments~\cite{BCS16}; evaluation-table
protocols such as FRI~\cite{BBHR18} and WHIR~\cite{ACFY24} motivate efficient
record hashing. The optimization applies when the existing leaf encoding
already consists of eight $n$-bit blocks and the entire record is disclosed.
For field-based data, eight blocks refers to the canonical bit encoding,
not necessarily to eight field elements.

If a prover constructs trees with $N_1,\ldots,N_r$ records, and the verifier
checks $u_j$ distinct records in tree $j$, the savings are
\begin{equation}
 \Delta_{\mathrm{prover}}=\sum_{j=1}^r N_j,
 \qquad
 \Delta_{\mathrm{verifier}}=\sum_{j=1}^r u_j
 \label{eq:snark}
\end{equation}
compression calls. Shared paths may be deduplicated as usual. The number
and sizes of roots, records, and authentication digests remain unchanged.
Their values change, however, so Fiat--Shamir challenges derived from the
commitments must be recomputed. Equal proof length does not mean identical
transcript bytes.

These gains concern the Merkle component. If it accounts for a fraction
$\alpha$ of prover time and a concrete implementation reduces that component
by a fraction $\beta$, the overall time becomes $(1-\alpha\beta)$ times
the original, assuming other costs stay fixed. A $20\%$ reduction in hash
calls does not by itself imply $\beta=0.2$, nor does it reduce polynomial
arithmetic, FFT, or sumcheck work.

Proposition~\ref{prop:binding} supplies consistency even for roots chosen by
a cheating prover. A particular compiler can require more, such as
extractability or a programmable random-oracle interface. Those requirements
must be checked separately; the binding argument alone is not a security
proof for every SNARK using the modified commitment. Challenge hashing
should also retain its own specified interface and separation.

If the original protocol already reveals the same records, the substitution
discloses no additional record contents. Packing previously separate query
locations into larger leaves is a different change: it may reveal more data
and affect query distributions, communication, or zero knowledge. The
same-size claim in this note does not include that change of granularity.

\subsection{Proving an opening inside a SNARK}
A second use is a SNARK that proves knowledge of a record and a valid Merkle
path. Its circuit evaluates the leaf hash and then the ordinary node hashes.
Let $g$ be the cost of each compression call and $g_{\oplus}$ the cost of
one $n$-bit XOR in the chosen circuit representation. With equal costs for
the specialized and baseline roles, the leaf-circuit saving is
\[
 4g-(3g+3g_{\oplus})=g-3g_{\oplus}.
\]
This may be favorable in a bit-oriented representation. In a prime-field
arithmetic circuit, XORs and bit decomposition can be expensive, so actual
constraint counts are needed. Replacing XOR by field addition defines a
different construction and does not automatically inherit the bitstring
security analysis.

\subsection{Other uses and implementation choices}
The same technique applies to authenticated file records, database rows,
fixed-size ledger entries, and outsourced data whenever complete records
are already the unit of authentication. The strongest total saving arises
when many leaves are hashed, as in initial commitment or re-commitment.
A verifier of one record still traverses the entire upper path. Tree
traversal, proof storage, and multiproof algorithms can operate over the
new leaf digests without changing their combinatorial structure.

A concrete evaluation should identify the encodings, security targets,
primitive implementation, and tree shape for both schemes. It should
include the cost of domain separation and report commitment throughput,
single-opening latency, batched verification, and memory use. For a SNARK,
it should also identify the fraction of total proving or verification time
spent on Merkle hashing. Section~\ref{sec:measurements} supplies measurements
of the Merkle component; an end-to-end performance claim still requires
the surrounding prover and verifier workload.

The main observation is simple: a complete record can be hashed more
cheaply while retaining the ordinary authentication tree above it. For the
specified eight-block format and baseline, this saves one compression call
per committed or verified record, with exactly the same opening size and
a direct reduction of position binding to complete-input collisions.

\paragraph{AI assistance.}
OpenAI Codex assisted with drafting and the analytical comparisons.

\begin{thebibliography}{FIPS202}
\bibitem[ACFY24]{ACFY24}
Gal Arnon, Alessandro Chiesa, Giacomo Fenzi, and Eylon Yogev.
\newblock \emph{WHIR: Reed--Solomon Proximity Testing with Super-Fast Verification}.
\newblock Cryptology ePrint Archive, Paper 2024/1586, 2024.
\newblock \url{https://eprint.iacr.org/2024/1586}.

\bibitem[BBHR18]{BBHR18}
Eli Ben-Sasson, Iddo Bentov, Yinon Horesh, and Michael Riabzev.
\newblock \emph{Fast Reed--Solomon Interactive Oracle Proofs of Proximity}.
\newblock ICALP 2018, LIPIcs 107, pages 14:1--14:17, 2018.
\newblock \url{https://doi.org/10.4230/LIPIcs.ICALP.2018.14}.

\bibitem[BCS16]{BCS16}
Eli Ben-Sasson, Alessandro Chiesa, and Nicholas Spooner.
\newblock \emph{Interactive Oracle Proofs}. TCC 2016-B, 2016.
\newblock \url{https://eprint.iacr.org/2016/116}.

\bibitem[DKMN21]{DKMN21}
Yevgeniy Dodis, Dmitry Khovratovich, Nicky Mouha, and Mridul Nandi.
\newblock \emph{T5: Hashing Five Inputs with Three Compression Calls}.
\newblock Cryptology ePrint Archive, Paper 2021/373, 2021.
\newblock \url{https://eprint.iacr.org/2021/373}.

\bibitem[BLAKE3]{BLAKE3}
The BLAKE3 team.
\newblock \emph{BLAKE3 specification}.
\newblock \url{https://github.com/BLAKE3-team/BLAKE3-specs/blob/master/blake3.tex}.

\bibitem[CN24]{CN24}
Debasmita Chakraborty and Mridul Nandi.
\newblock \emph{Lower Bound on Number of Compression Calls of a
Collision-Resistance Preserving Hash}.
\newblock Cryptology ePrint Archive, Paper 2024/1095, 2024.
\newblock \url{https://eprint.iacr.org/2024/1095}.

\bibitem[DDN22]{DDN22}
Chandranan Dhar, Yevgeniy Dodis, and Mridul Nandi.
\newblock \emph{Revisiting Collision and Local Opening Analysis of ABR Hash}.
\newblock ITC 2022, LIPIcs 230, pages 11:1--11:22, 2022.
\newblock \url{https://doi.org/10.4230/LIPIcs.ITC.2022.11}.

\bibitem[FIPS202]{FIPS202}
National Institute of Standards and Technology.
\newblock \emph{SHA-3 Standard: Permutation-Based Hash and Extendable-Output
Functions}. FIPS PUB 202, August 2015.
\newblock \url{https://doi.org/10.6028/NIST.FIPS.202}.

\bibitem[SLZ+25]{SPONGEDM}
Siwei Sun, Shun Li, Zhiyu Zhang, Charlotte Lefevre, Bart Mennink,
Zhen Qin, and Dengguo Feng.
\newblock \emph{Permutation-Based Hashing With Stronger (Second) Preimage
Resistance}.
\newblock Cryptology ePrint Archive, Paper 2025/963, revised June 2026.
\newblock \url{https://eprint.iacr.org/2025/963}.
\end{thebibliography}
\end{document}
